
\documentclass[12pt,a4paper]{report}

% Packages
% Page geometry comes from styles.PAGE: header and footer heights are reserved
% explicitly (includehead/includefoot), so running heads/feet never overlap
% the text block.
\usepackage[a4paper, left=2.54cm, right=2.54cm, top=2.54cm, bottom=2.54cm, headheight=15.0pt, headsep=14.0pt, footskip=42.0pt, includehead, includefoot]{geometry}
\usepackage{graphicx}
\usepackage{amsmath}
\usepackage{amssymb}
\usepackage{booktabs}
\usepackage{array}
\usepackage{tabularx}
\usepackage{float}
\usepackage{fancyhdr}
\usepackage[hidelinks]{hyperref}
\usepackage{xcolor}
\usepackage{setspace}
\usepackage{enumitem}
\usepackage{caption}

\captionsetup{
    labelfont=bf,
    justification=raggedright,
    singlelinecheck=false,
    format=plain
}
\usepackage{subcaption}
\usepackage{multirow}
\usepackage{colortbl}
\usepackage{longtable}
\setlength{\LTleft}{\fill}
\setlength{\LTright}{\fill}
\usepackage{titlesec}
\usepackage{titletoc}
\usepackage{lastpage}
\usepackage{makecell}
\usepackage{etoolbox}
\usepackage{needspace}

\numberwithin{table}{chapter}
\numberwithin{figure}{chapter}
% Table layout and spacing: consistent padding, row height, and longtable pre/post skips
% (values from styles.TABLE -- single source of truth)
\setlength{\tabcolsep}{6.0pt}
\renewcommand{\arraystretch}{1.12}
\setlength{\LTpre}{0.0pt}
\setlength{\LTpost}{6.0pt}
\setlength{\arrayrulewidth}{0.5pt}
\setlength{\extrarowheight}{0.6pt}
\setcounter{LTchunksize}{1000}

% A partial rule (\cline) is set as its own table line followed by
% \vskip-\arrayrulewidth; that glue is a legal page break, so a page could
% break right after a \cline -- in the middle of a \multirow girder group,
% pushing the group label into the footer.  Forbid a break at that point.
\makeatletter
\patchcmd{\@cline}{\noalign{\vskip-\arrayrulewidth}}{\noalign{\nobreak\vskip-\arrayrulewidth}}{}{\typeout{OsdagBridge: cline patch not applied}}
\makeatother

% Prevent (float) tables from starting at the very bottom of a page:
% if fewer than N baseline-skips remain, break to the next page first.
% NOTE: deliberately NOT applied to longtable.  A \needspace break issued
% right before a longtable is processed by longtable's own output routine,
% which then prints the "continued" foot alone on the old page and both the
% continued and the first heading on the new page.  longtable already moves
% its heading + first row to the next page when they do not fit.
\BeforeBeginEnvironment{table}{\needspace{5\baselineskip}}

% Report colours (styles.COLORS / styles.LATEX_COLOR_NAMES)
\definecolor{osdagGreen}{HTML}{91B014}
\definecolor{osdagSecondary}{HTML}{4D5B23}
\definecolor{osdagAccent}{HTML}{1F77B4}
\definecolor{tableHeaderBg}{HTML}{EEF3DC}
\definecolor{passGreen}{HTML}{2E8B57}
\definecolor{failRed}{HTML}{C62828}
\definecolor{mutedGrey}{HTML}{6B6B6B}

% Header / footer (styles.HEADER_FOOTER).  The optional "* Software default
% value" note is part of the footer block, inside the space reserved by
% \footskip, instead of being pulled up into the text area with a negative
% \vspace (which made it overlap the last table row of the page).
\fancypagestyle{main}{
  \fancyhf{}
  \fancyhead[L]{OsdagBridge Demonstration Case (30 m span) $|$ DEMO-30}
  \fancyhead[R]{2026-11-23 $|$ Rev 0}
  \fancyfoot[L]{Osdag $|$ FOSSEE $|$ Indian Institute of Technology Bombay}
  \fancyfoot[R]{Page \thepage\ of \pageref{LastPage}}
  \renewcommand{\headrule}{{\color{osdagGreen}\hrule width\headwidth height 1.0pt}}
  \renewcommand{\footrule}{%
    \ifbool{hasSDonPage}{%
      \hbox to \headwidth{\footnotesize\itshape\textcolor{black}{* Software default value}\hfil}%
      \vskip 2.0pt%
    }{}%
    {\color{osdagGreen}\hrule width\headwidth height 1.0pt}%
    \vskip 4.0pt%
  }
}
\fancypagestyle{plain}{
  \fancyhf{}
  \fancyhead[L]{OsdagBridge Demonstration Case (30 m span) $|$ DEMO-30}
  \fancyhead[R]{2026-11-23 $|$ Rev 0}
  \fancyfoot[L]{Osdag $|$ FOSSEE $|$ Indian Institute of Technology Bombay}
  \fancyfoot[R]{Page \thepage\ of \pageref{LastPage}}
  \renewcommand{\headrule}{{\color{osdagGreen}\hrule width\headwidth height 1.0pt}}
  \renewcommand{\footrule}{%
    \ifbool{hasSDonPage}{%
      \hbox to \headwidth{\footnotesize\itshape\textcolor{black}{* Software default value}\hfil}%
      \vskip 2.0pt%
    }{}%
    {\color{osdagGreen}\hrule width\headwidth height 1.0pt}%
    \vskip 4.0pt%
  }
}
\fancypagestyle{firstpage}{
  \fancyhf{}
  \renewcommand{\headrulewidth}{0pt}
  \fancyfoot[L]{Osdag $|$ FOSSEE $|$ Indian Institute of Technology Bombay}
  % No page number printed on the title page (it is page 1; numbering
  % continues at 2 after the titlepage environment, see title_page()).
  \fancyfoot[R]{}
  \renewcommand{\footrule}{{\color{osdagGreen}\hrule width\headwidth height 1.0pt}\vskip 4.0pt}
}
\pagestyle{main}
\setstretch{1.15}
\setlength{\emergencystretch}{3em}

% Custom Commands
\newcommand{\placeholder}[1]{\textit{\textless #1\textgreater}}
\newcommand{\todo}[1]{\colorbox{yellow}{TODO: #1}}
\newcolumntype{L}[1]{>{\raggedright\arraybackslash}p{#1}}
\newcolumntype{C}[1]{>{\centering\arraybackslash}p{#1}}
\newcolumntype{R}[1]{>{\raggedleft\arraybackslash}p{#1}}

% Software-default asterisk
\newcommand{\sdstar}{\textsuperscript{*}}
\newbool{hasSDonPage}
\boolfalse{hasSDonPage}
\newcommand{\markSD}{\global\booltrue{hasSDonPage}}
\renewcommand{\sdstar}{\textsuperscript{*}\markSD{}}
\AddToHook{shipout/before}{\global\boolfalse{hasSDonPage}}

\title{\Large\textbf{OsdagBridge} \\ \normalsize Open Source Software for Steel Girder Bridge Design \\ \vspace{2cm} \large Design Report}
\author{}
\date{}

\begin{document}


\begin{titlepage}
\thispagestyle{firstpage}
\centering
\vspace*{1.5cm}
\noindent
\begin{minipage}[c]{0.6\textwidth}
\raggedright
\includegraphics[width=\linewidth,keepaspectratio]{assets/osdag_logo.png}
\end{minipage}%
\hfill
\begin{minipage}[c]{0.35\textwidth}
\raggedleft

\end{minipage}
\\[1cm]

{\Huge \textbf{OsdagBridge}}\\[0.3cm]
{\large Open Source Software for Steel Girder Bridge Design}\\[1.5cm]
{\Large Design Report}\\[1.5cm]
\begin{tabular}{|L{4cm}|L{10cm}|}
\hline
\textbf{Project Name} & OsdagBridge Demonstration Case (30 m span) \\
\hline
\textbf{Project Location} & Mumbai(Bombay) (Santa Cruz), Maharashtra \\
\hline
\textbf{Author / Designer} & Report generator QA \\
\hline
\textbf{Reviewer} &  \\
\hline
\textbf{Organization} &  \\
\hline
\textbf{Client Name and Organization} &  \\
\hline
\textbf{Job Number} & DEMO-30 \\
\hline
\textbf{Date} & 2026-11-23 \\
\hline
\textbf{Report Version} & Rev 0 \\
\hline
\end{tabular}
\end{titlepage}
% The report class resets the page counter to 1 after titlepage; continue at
% 2 so that "Page X of Y" equals the PDF page number (title page = page 1,
% printed without a number) and Y equals the PDF page count.
\setcounter{page}{2}


% Chapter / TOC Formatting
\titleformat{\chapter}[block]
  {\normalfont\Large\bfseries\centering}{\thechapter}{1em}{}
\titlespacing*{\chapter}{0pt}{0pt}{10pt}
\setcounter{tocdepth}{2}

% TOC styling using titletoc
\titlecontents{chapter}[1.5em]
  {\normalfont\vspace{2pt}}
  {\contentslabel{1.5em}}
  {\hspace*{-1.5em}}
  {\hfill\contentspage}

\titlecontents{section}[3.8em]
  {\normalfont}
  {\contentslabel{2.3em}}
  {\hspace*{-2.3em}}
  {\hfill\contentspage}

\titlecontents{subsection}[7.0em]
  {\normalfont}
  {\contentslabel{3.2em}}
  {\hspace*{-3.2em}}
  {\hfill\contentspage}

\newpage
\renewcommand{\contentsname}{\centering\Large\bfseries Table of Contents}
\tableofcontents


\newpage
{\centering\Large\bfseries Executive Summary\par}
\addcontentsline{toc}{chapter}{Executive Summary}
\vspace{0.8em}

This section provides a concise summary of the bridge design, key inputs, governing loads, and final design outcomes.

\section*{Project Overview}
\addcontentsline{toc}{section}{Project Overview}
\label{sec:project-overview}


\begin{tabular}{|L{5.5cm}|L{8.5cm}|}
\hline
\textbf{Bridge Type} & Highway Bridge \\
\hline
\textbf{Design Standard} & IRC 5, IRC 6, IRC 22, IRC 24, IS 800 \\
\hline
\textbf{Span} & 30 m \\
\hline
\textbf{Carriageway Width} & 7.5 m \\
\hline
\textbf{No. of Girders} & 4 \\
\hline
\textbf{Girder Spacing} & 2.1 m \\
\hline
\textbf{Deck Thickness} & 250.0 mm \\
\hline
\textbf{Overall Design Status} & \textcolor{red}{\textbf{FAIL}} --- Fatigue Normal Stress (G3), UR 1.58; Int.Stiff:\allowbreak{} Leg \ensuremath{\leq} H\_\allowbreak{}limit (G1), UR 1.48 \\
\hline
\textbf{Governing Check} & Fatigue Normal Stress (G3) --- Steel Plate Girders \\
\hline
\textbf{Overall Utilization Ratio (max)} & \textcolor{red}{1.58} (Steel Plate Girders) \\
\hline
\end{tabular}


\begin{figure}[H]
\centering
\includegraphics[width=0.9\textwidth]{C:/Users/udayk/AppData/Local/Temp/tmpymqxq2pd/images/girder_top.png}
\caption*{Figure 1 -- Overall Bridge Plan}
\end{figure}

\newpage

\begin{figure}[H]
\centering
\includegraphics[width=0.9\textwidth]{C:/Users/udayk/AppData/Local/Temp/tmpymqxq2pd/images/cross_section.png}
\caption*{Figure 2 -- Typical Cross-Section (with girder, deck, barriers, footpath)}
\end{figure}

\begin{figure}[H]
\centering
\includegraphics[width=0.9\textwidth]{C:/Users/udayk/AppData/Local/Temp/tmpymqxq2pd/images/final_geometry.png}
\caption*{Figure 3 -- 3D View of Bridge Superstructure}
\end{figure}

\par\noindent\begin{minipage}{\textwidth}
\noindent\textbf{Table 1 -- Final Bridge Geometry (after optimization)}

\vspace{0.4em}
\noindent
\begin{tabular}{|C{2.58cm}|C{2.76cm}|C{2.76cm}|C{2.76cm}|C{2.76cm}|}
\hline
  \textbf{} &
  \textbf{G1} &
  \textbf{G2} &
  \textbf{G3} &
  \textbf{G4} \\
\hline
Member ID & G1M1 & G2M1 & G3M1 & G4M1 \\
\hline
Section Designation & 1670 x 510 x 22 x 510 x 22 & 1670 x 510 x 22 x 510 x 22 & 1670 x 510 x 22 x 510 x 22 & 1670 x 510 x 22 x 510 x 22 \\
\hline
Governing Check & Fatigue Normal Stress & Fatigue Normal Stress & Fatigue Normal Stress & Fatigue Normal Stress \\
\hline
Utilization Ratio (max) & \textcolor{red}{1.54} (\textcolor{red}{FAIL}) & \textcolor{red}{1.54} (\textcolor{red}{FAIL}) & \textcolor{red}{1.58} (\textcolor{red}{FAIL}) & \textcolor{red}{1.58} (\textcolor{red}{FAIL}) \\
\hline
\end{tabular}
\end{minipage}

\vspace{0.4em}
\noindent\textit{Note: Utilization ratio (UR) = demand / capacity. UR $\leq 1.0$ passes, UR $> 1.0$ fails. Girder columns show each girder's governing check as labelled by the design module.}

\vspace{1em}

\section*{Key Design Outcomes Summary}
\addcontentsline{toc}{section}{Key Design Outcomes Summary}
\label{sec:key-outcomes}

\noindent\begin{tabular}{|L{3.01cm}|C{2.12cm}|L{5.32cm}|C{1.24cm}|C{1.95cm}|}
\hline
\rowcolor{tableHeaderBg}\textbf{Component} & \textbf{Status} & \textbf{Governing check} & \textbf{Max UR} & \textbf{Checks with data} \\
\hline
Steel Plate Girders & \textcolor{red}{\textbf{FAIL}} & Fatigue Normal Stress (G3) & \textcolor{red}{1.58} & 20 of 20 \\
\hline
Concrete Deck Slab & \textbf{PASS} & Cantilever overhang & 0.94 & 6 of 6 \\
\hline
Cross Bracing & \textbf{PASS} & G1-G2 diagonal (C) & 0.74 & 12 of 12 \\
\hline
End Diaphragms & \textbf{NOT AVAILABLE} & The analysis returned no end-diaphragm member forces, so no end-diaphragm Demand/Capacity checks were performed. & --- & 0 of 0 \\
\hline
\end{tabular}

\vspace{0.4em}
\noindent\textit{Note: A component passes only when every one of its checks has valid Demand / Capacity data and UR $\leq 1.0$; a component without results is ``NOT AVAILABLE'' and is never counted as passing. Details: Section 5.5.}

\section*{Data Consistency Notes}
\addcontentsline{toc}{section}{Data Consistency Notes}
The following inconsistencies were found in the analysis / design output. They are reported here as found; the report does not change engineering results.
\begin{itemize}
\item The girder design results contain 6 longitudinal members (G1-G6) although the bridge has 4 girders. The grillage model adds a longitudinal line at each deck edge, and the design module designs and labels these edge lines as girders as well. Results are reported with the labels used by the design module, so these labels do not necessarily denote the girders numbered in Chapter 2. This is an issue of the analysis/design module (not changed by the report) and requires engineering verification.
\item Cross-bracing spacing:\allowbreak{} the input lists 15 m (Table 2.9), but the girder design checks and the cross-bracing force calculation used 3 m. The design modules read a spacing value that is not set by the input, so their default was applied. Requires engineering verification.
\item Girder section properties:\allowbreak{} the grillage analysis model uses the stored section area 349.2 cm2, whereas the final designed section (1670 x 510 x 22 x 510 x 22) has 387.0 cm2. Member forces were therefore computed with properties that differ from the checked section. Requires engineering verification.
\item Deflection:\allowbreak{} the per-girder load-effect results used in Table 4.3 give 19.385 mm (live) and 51.225 mm (total), whereas the IRC 22 deflection checks of the girder design module (Section 5.1) use 9.631 mm and 22.871 mm. The two calculations of the application disagree, so the deflection status differs between Table 4.3 and Section 5. Requires engineering verification.
\end{itemize}

\section*{Design Assumptions and Limitations}
\addcontentsline{toc}{section}{Design Assumptions and Limitations}
\label{sec:assumptions}

\begin{itemize}
\item Additional inputs not provided by the user were assumed by software per IRC/IS code defaults or practical consideration.
\item Grillage analysis was performed using OSPGrillage assuming simply supported I-girders.
\item Substructure and foundation design are not included in this report.
\item Splice connections and bearings are not designed in this version.
\end{itemize}

% Restore numbered chapter format
\titleformat{\chapter}[block]{\normalfont\Large\bfseries\centering}{\thechapter}{1em}{}
\titlespacing*{\chapter}{0pt}{-30pt}{10pt}


\chapter{Project Information}

This section records all project metadata as entered by the designer.

\section{Project and Design Team Details}
\label{sec:project-details}

\begin{tabular}{|L{5.5cm}|L{8.5cm}|}
\hline
\textbf{Project Name} & OsdagBridge Demonstration Case (30 m span) \\
\hline
\textbf{Project Location} & Mumbai(Bombay) (Santa Cruz), Maharashtra \\
\hline
\textbf{Designer} & Report generator QA \\
\hline
\textbf{Reviewer} &  \\
\hline
\textbf{Organization} &  \\
\hline
\textbf{Client} &  \\
\hline
\textbf{Software Version} & OsdagBridge \\
\hline
\end{tabular}


\section{Applicable Codes and Standards}
\label{sec:codes}

\begin{itemize}
\item Indian Roads Congress (IRC) 5: General Features of Design
\item Indian Roads Congress (IRC) 6: Loads and Load Combinations
\item Indian Roads Congress (IRC) 22: Composite Construction (Limit State Design)
\item Indian Roads Congress (IRC) 24: Steel Road Bridges (Limit State Method)
\item Indian Roads Congress (IRC) 112: Concrete Road Bridges (deck design)
\item Indian Roads Congress Special Publication (IRC SP) 114: Seismic Design of Road Bridges
\item Indian Standard (IS) 800: General Construction in Steel
\item Indian Standard (IS) 2062: Hot Rolled Structural Steel Specification
\item Indian Standard (IS) 6006: Steel Bearings
\end{itemize}


\chapter{Input Parameters}

\setlength{\abovecaptionskip}{2pt}
\setlength{\belowcaptionskip}{2pt}

This section documents all inputs provided to OsdagBridge. User-provided inputs are clearly distinguished from software-assumed defaults. Where the user did not supply a value, the software has applied the IRC/IS code default or an empirical guideline; these are annotated with an asterisk (\sdstar{}).

\section{Basic Inputs (User-Defined)}
\label{sec:basic-inputs}

\noindent\textit{Note: These inputs are mandatory and were provided by the user.}

\begin{table}[H]
\caption{\textbf{Project Location}}
\label{subsec:project-location}
\begin{tabular}{|L{5.5cm}|L{8.5cm}|}
\hline
\textbf{parameter} & \textbf{value} \\
\hline
\textnormal{Project Location} & Mumbai(Bombay) (Santa Cruz), Maharashtra \\
\hline
\textnormal{Latitude / Longitude} & 19.076, 72.8777 \\
\hline
\textnormal{Seismic Zone (IRC 6)} & III \\
\hline
\textnormal{Basic Wind Speed (IRC 6)} & 44 m/s \\
\hline
\textnormal{Shade Temp. Max / Min (IRC 6)} & 42.2 °C / 7.4 °C \\
\hline
\end{tabular}
\vspace{0.4cm}
\end{table}

\begin{table}[H]
\caption{\textbf{Bridge Geometry}}
\label{subsec:bridge-geometry}
\begin{tabular}{|L{5.5cm}|L{8.5cm}|}
\hline
\textbf{parameter} & \textbf{value} \\
\hline
\textnormal{Type of Structure} & Highway Bridge \\
\hline
\textnormal{Span (m)} & 30 m \\
\hline
\textnormal{Carriageway Width (m)} & 7.5 m \\
\hline
\textnormal{Include Median} & No \\
\hline
\textnormal{Footpath} & None \\
\hline
\textnormal{Skew Angle (degrees)} & 0° (IRC 24 Cl. 504.8 limit: $\pm$15°) \\
\hline
\end{tabular}
\end{table}
\vspace{0.4cm}

\begin{table}[H]
\caption{\textbf{Material Selection}}
\label{subsec:material}
\begin{tabular}{|L{5.5cm}|L{8.5cm}|}
\hline
\textbf{parameter} & \textbf{value} \\
\hline
\textnormal{Girder Steel Grade (IS 2062)} & E 350A \\
\hline
\textnormal{Cross Bracing Steel Grade} & E 350A \\
\hline
\textnormal{End Diaphragm Steel Grade} & E 350A \\
\hline
\textnormal{Concrete Deck Grade (IRC 22)} & M40 \\
\hline
\end{tabular}
\end{table}
\vspace{0.4cm}

\newpage
\section{Additional Inputs}
\label{sec:additional-inputs}

Where the user has modified additional inputs, those values are reported here. Where no modification was made, the software default is shown.

\vspace{0.8cm}

\begin{longtable}{|L{5.31cm}|p{9.66cm}|}%
\caption{\textbf{Typical Section Details}}\\%
\hline%
\rowcolor{tableHeaderBg}\hspace{0pt}\textbf{parameter} & \hspace{0pt}\textbf{value} \\%
\hline%
\endfirsthead%
\caption[]{\textbf{Typical Section Details} (continued)}\\%
\hline%
\rowcolor{tableHeaderBg}\hspace{0pt}\textbf{parameter} & \hspace{0pt}\textbf{value} \\%
\hline%
\endhead%
\multicolumn{2}{r}{{\footnotesize\itshape Continued on next page}} \\%
\endfoot%
\hline%
\endlastfoot%
\textnormal{Overall Bridge Width (m)} & 8.4 \\*[6pt]
\hline
\textnormal{No. of Girders} & 4 \\*[6pt]
\hline
\textnormal{Girder Spacing (m)} & 2.1 m \\*[6pt]
\hline
\textnormal{Deck Overhang Width (m)} & 1.05 m \\*[6pt]
\hline
\textnormal{Deck Thickness (mm)} & 250.0 mm \\*[6pt]
\hline
\textnormal{Footpath Width (m)} & 0.0 m (IRC 5 Cl. 104.3.6 min: 1.5 m) \\*[6pt]
\hline
\textnormal{No. of Traffic Lanes} & 2 [IRC 6 Table 6 (carriageway width)] \\*[6pt]%
\end{longtable}

\vspace{0.8em}

\begin{longtable}{|L{5.31cm}|p{9.66cm}|}%
\caption{\textbf{Components Details}}\\%
\hline%
\rowcolor{tableHeaderBg}\hspace{0pt}\textbf{parameter} & \hspace{0pt}\textbf{value} \\%
\hline%
\endfirsthead%
\caption[]{\textbf{Components Details} (continued)}\\%
\hline%
\rowcolor{tableHeaderBg}\hspace{0pt}\textbf{parameter} & \hspace{0pt}\textbf{value} \\%
\hline%
\endhead%
\multicolumn{2}{r}{{\footnotesize\itshape Continued on next page}} \\%
\endfoot%
\hline%
\endlastfoot%
\textnormal{Crash Barrier Type} & IRC 5 - RCC Crash Barrier \\*[6pt]
\hline
\textnormal{Crash Barrier Load (kN/m)} & 7.56 \\*[6pt]
\hline

\textnormal{Railing Type} & IRC 5 - RCC Railing \\*[6pt]
\hline
\textnormal{Railing Load (kN/m)} & 6.806 \\*[6pt]
\hline
\textnormal{Wearing Course Material} & Concrete \\*[6pt]
\hline
\textnormal{Wearing Course Thickness (mm)} & 50.0 mm \\*[6pt]%
\end{longtable}


\newpage

\vspace{0.4em}
\noindent
            
\vspace{4pt}
\begin{longtable}{|L{1.95cm}|L{1.59cm}|p{3.36cm}|p{3.36cm}|p{3.36cm}|}%
\caption{\textbf{Girder General Information}}\\%
\hline%
\rowcolor{tableHeaderBg}\hspace{0pt}\textbf{Girder} & \hspace{0pt}\textbf{Member ID} & \hspace{0pt}\textbf{Design Mode} & \hspace{0pt}\textbf{Girder Type} & \hspace{0pt}\textbf{Girder Symmetry} \\%
\hline%
\endfirsthead%
\caption[]{\textbf{Girder General Information} (continued)}\\%
\hline%
\rowcolor{tableHeaderBg}\hspace{0pt}\textbf{Girder} & \hspace{0pt}\textbf{Member ID} & \hspace{0pt}\textbf{Design Mode} & \hspace{0pt}\textbf{Girder Type} & \hspace{0pt}\textbf{Girder Symmetry} \\%
\hline%
\endhead%
\multicolumn{5}{r}{{\footnotesize\itshape Continued on next page}} \\%
\endfoot%
\hline%
\endlastfoot%
G1 & G1M1 & Optimized & Welded & Girder Symmetric \\*[8pt]
\hline
G2 & G2M1 & Optimized & Welded & Girder Symmetric \\*[8pt]
\hline
G3 & G3M1 & Optimized & Welded & Girder Symmetric \\*[8pt]
\hline
G4 & G4M1 & Optimized & Welded & Girder Symmetric \\*[8pt]%
\end{longtable}

\vspace{0.6em}

\vspace{4pt}
\begin{longtable}{|L{1.58cm}|L{2.02cm}|L{1.58cm}|p{4.22cm}|p{4.22cm}|}%
\caption{\textbf{Girder Section Dimensions}}\\%
\hline%
\rowcolor{tableHeaderBg}\hspace{0pt}\textbf{Girder} & \hspace{0pt}\textbf{Total Depth, D (mm)} & \hspace{0pt}\textbf{Web, $t_w$ (mm)} & \hspace{0pt}\textbf{Top Flange (b\textsubscript{tf}, t\textsubscript{tf}) mm} & \hspace{0pt}\textbf{Bottom Flange (b\textsubscript{bf}, t\textsubscript{bf}) mm} \\%
\hline%
\endfirsthead%
\caption[]{\textbf{Girder Section Dimensions} (continued)}\\%
\hline%
\rowcolor{tableHeaderBg}\hspace{0pt}\textbf{Girder} & \hspace{0pt}\textbf{Total Depth, D (mm)} & \hspace{0pt}\textbf{Web, $t_w$ (mm)} & \hspace{0pt}\textbf{Top Flange (b\textsubscript{tf}, t\textsubscript{tf}) mm} & \hspace{0pt}\textbf{Bottom Flange (b\textsubscript{bf}, t\textsubscript{bf}) mm} \\%
\hline%
\endhead%
\multicolumn{5}{r}{{\footnotesize\itshape Continued on next page}} \\%
\endfoot%
\hline%
\endlastfoot%
G1 & 1670 mm & 10 mm & 510 mm, 22 mm & 510 mm, 22 mm \\*[8pt]
\hline
G2 & 1670 mm & 10 mm & 510 mm, 22 mm & 510 mm, 22 mm \\*[8pt]
\hline
G3 & 1670 mm & 10 mm & 510 mm, 22 mm & 510 mm, 22 mm \\*[8pt]
\hline
G4 & 1670 mm & 10 mm & 510 mm, 22 mm & 510 mm, 22 mm \\*[8pt]%
\end{longtable}

\vspace{0.6em}

\vspace{4pt}
\begin{longtable}{|L{1.38cm}|p{2.16cm}|p{2.16cm}|p{2.95cm}|p{2.36cm}|p{2.16cm}|}%
\caption{\textbf{Girder Restraint and Stiffener Details}}\\%
\hline%
\rowcolor{tableHeaderBg}\hspace{0pt}\textbf{Girder} & \hspace{0pt}\textbf{Torsional / Warping Restraint} & \hspace{0pt}\textbf{Web Philosophy} & \hspace{0pt}\textbf{Intermediate Stiffeners} & \hspace{0pt}\textbf{Longitudinal Stiffeners} & \hspace{0pt}\textbf{Bearing Stiffener} \\%
\hline%
\endfirsthead%
\caption[]{\textbf{Girder Restraint and Stiffener Details} (continued)}\\%
\hline%
\rowcolor{tableHeaderBg}\hspace{0pt}\textbf{Girder} & \hspace{0pt}\textbf{Torsional / Warping Restraint} & \hspace{0pt}\textbf{Web Philosophy} & \hspace{0pt}\textbf{Intermediate Stiffeners} & \hspace{0pt}\textbf{Longitudinal Stiffeners} & \hspace{0pt}\textbf{Bearing Stiffener} \\%
\hline%
\endhead%
\multicolumn{6}{r}{{\footnotesize\itshape Continued on next page}} \\%
\endfoot%
\hline%
\endlastfoot%
G1 & Fully Restrained, Both Flanges Restrained & Thin Web with ITS & Yes; Spacing: 2440 mm; Thickness: 10 mm & No & No.: 4; Spacing: 100 mm; Thickness: 10 mm \\*[8pt]
\hline
G2 & Fully Restrained, Both Flanges Restrained & Thin Web with ITS & Yes; Spacing: 2440 mm; Thickness: 10 mm & No & No.: 4; Spacing: 100 mm; Thickness: 10 mm \\*[8pt]
\hline
G3 & Fully Restrained, Both Flanges Restrained & Thin Web with ITS & Yes; Spacing: 2440 mm; Thickness: 10 mm & No & No.: 4; Spacing: 100 mm; Thickness: 10 mm \\*[8pt]
\hline
G4 & Fully Restrained, Both Flanges Restrained & Thin Web with ITS & Yes; Spacing: 2440 mm; Thickness: 10 mm & No & No.: 4; Spacing: 100 mm; Thickness: 10 mm \\*[8pt]%
\end{longtable}



\newpage

\vspace{0.4em}

\begin{longtable}{|L{2.15cm}|L{2.15cm}|L{2.93cm}|L{2.44cm}|C{1.76cm}|C{1.76cm}|}%
\caption{\textbf{Member Properties: Cross Bracing Details}}\\%
\hline%
\rowcolor{tableHeaderBg}\hspace{0pt}\textbf{Location} & \hspace{0pt}\textbf{Member IDs} & \hspace{0pt}\textbf{Type of Bracing} & \hspace{0pt}\textbf{Bracing Section} & \hspace{0pt}\textbf{Spacing (m)} \\%
\hline%
\endfirsthead%
\caption[]{\textbf{Member Properties: Cross Bracing Details} (continued)}\\%
\hline%
\rowcolor{tableHeaderBg}\hspace{0pt}\textbf{Location} & \hspace{0pt}\textbf{Member IDs} & \hspace{0pt}\textbf{Type of Bracing} & \hspace{0pt}\textbf{Bracing Section} & \hspace{0pt}\textbf{Spacing (m)} \\%
\hline%
\endhead%
\multicolumn{6}{r}{{\footnotesize\itshape Continued on next page}} \\%
\endfoot%
\hline%
\endlastfoot%
G1 to G2 & B1M1 to B1M1 & X-Bracing & IS 100 x 100 x 10 & 15.0 m \\*[6pt]
\hline
G2 to G3 & B2M1 to B2M1 & X-Bracing & IS 100 x 100 x 10 & 15.0 m \\*[6pt]
\hline
G3 to G4 & B3M1 to B3M1 & X-Bracing & IS 100 x 100 x 10 & 15.0 m \\*[6pt]%
\end{longtable}

\vspace{0.4em}
\noindent

\begin{longtable}{|L{2.15cm}|L{2.15cm}|L{2.93cm}|L{2.44cm}|C{1.76cm}|C{1.76cm}|}%
\caption{\textbf{Member Properties: End Diaphragm Details}}\\%
\hline%
\rowcolor{tableHeaderBg}\hspace{0pt}\textbf{Location} & \hspace{0pt}\textbf{Member IDs} & \hspace{0pt}\textbf{Type of Bracing} & \hspace{0pt}\textbf{Bracing Section} \\%
\hline%
\endfirsthead%
\caption[]{\textbf{Member Properties: End Diaphragm Details} (continued)}\\%
\hline%
\rowcolor{tableHeaderBg}\hspace{0pt}\textbf{Location} & \hspace{0pt}\textbf{Member IDs} & \hspace{0pt}\textbf{Type of Bracing} & \hspace{0pt}\textbf{Bracing Section} \\%
\hline%
\endhead%
\multicolumn{6}{r}{{\footnotesize\itshape Continued on next page}} \\%
\endfoot%
\hline%
\endlastfoot%
G1 to G2 & E1M1 to E1M2 & Cross Bracing & IS 100 x 100 x 10 \\*[6pt]
\hline
G2 to G3 & E2M1 to E2M2 & Cross Bracing & IS 100 x 100 x 10 \\*[6pt]
\hline
G3 to G4 & E3M1 to E3M2 & Cross Bracing & IS 100 x 100 x 10 \\*[6pt]%
\end{longtable}

\noindent\textit{Note: Cross-bracing spacing:\allowbreak{} the input lists 15 m (Table 2.9), but the girder design checks and the cross-bracing force calculation used 3 m. The design modules read a spacing value that is not set by the input, so their default was applied. Requires engineering verification.}



\label{subsec:shear-connectors}

\vspace{2.2em}

\vspace{0.4em}
\begin{longtable}{|L{5.31cm}|p{9.66cm}|}%
\caption{\textbf{Shear Connector Details}}\\%
\hline%
\rowcolor{tableHeaderBg}\hspace{0pt}\textbf{parameter} & \hspace{0pt}\textbf{value} \\%
\hline%
\endfirsthead%
\caption[]{\textbf{Shear Connector Details} (continued)}\\%
\hline%
\rowcolor{tableHeaderBg}\hspace{0pt}\textbf{parameter} & \hspace{0pt}\textbf{value} \\%
\hline%
\endhead%
\multicolumn{2}{r}{{\footnotesize\itshape Continued on next page}} \\%
\endfoot%
\hline%
\endlastfoot%
\textnormal{Stud Diameter (mm)} & 20.0 mm \\*[6pt]
\hline
\textnormal{Stud Height (mm)} & 100.0 mm \\*[6pt]
\hline
\textnormal{Stud $f_y$ (MPa)} & 385.0 MPa \\*[6pt]
\hline
\textnormal{Stud $f_u$ (MPa)} & 495.0 MPa \\*[6pt]
\hline
\textnormal{No. of Studs per Section} & 2 \\*[6pt]%
\end{longtable}



\label{subsec:safety-factors}

\vspace{2.2em}

\vspace{0.3em}
\noindent\textit{Note: All values are per IRC 22 Table 1 unless user-modified.}

\vspace{0.4em}
\begin{longtable}{|L{5.31cm}|p{9.66cm}|}%
\caption{\textbf{Partial Safety Factors}}\\%
\hline%
\rowcolor{tableHeaderBg}\hspace{0pt}\textbf{parameter} & \hspace{0pt}\textbf{value} \\%
\hline%
\endfirsthead%
\caption[]{\textbf{Partial Safety Factors} (continued)}\\%
\hline%
\rowcolor{tableHeaderBg}\hspace{0pt}\textbf{parameter} & \hspace{0pt}\textbf{value} \\%
\hline%
\endhead%
\multicolumn{2}{r}{{\footnotesize\itshape Continued on next page}} \\%
\endfoot%
\hline%
\endlastfoot%
\textnormal{$\gamma_{M0}$ (Yielding / Buckling)} & 1.1 \\*[6pt]
\hline
\textnormal{$\gamma_{M1}$ (Ultimate Stress)} & 1.25 \\*[6pt]
\hline
\textnormal{$\gamma_C$ (Concrete, Basic)} & 1.5 \\*[6pt]
\hline
\textnormal{$\gamma_s$ (Reinforcement)} & 1.15 \\*[6pt]
\hline
\textnormal{$\gamma_v$ (Shear Connectors)} & 1.25 \\*[6pt]
\hline
\textnormal{$\gamma_{fft}$ (Fatigue Load)} & 1.0 \\*[6pt]
\hline
\textnormal{$\gamma_{Mft}$ (Fatigue Strength)} & 1.35 \\*[6pt]%
\end{longtable}


\chapter{Loads and Load Combinations}

This section summarizes all loads applied to the bridge and the load combinations considered for analysis and design.

\vspace{1em}
\begin{longtable}{|L{5.31cm}|p{9.66cm}|}%
\caption{\textbf{Dead Load -- Self Weight}}\\%
\hline%
\rowcolor{tableHeaderBg}\hspace{0pt}\textbf{parameter} & \hspace{0pt}\textbf{value} \\%
\hline%
\endfirsthead%
\caption[]{\textbf{Dead Load -- Self Weight} (continued)}\\%
\hline%
\rowcolor{tableHeaderBg}\hspace{0pt}\textbf{parameter} & \hspace{0pt}\textbf{value} \\%
\hline%
\endhead%
\multicolumn{2}{r}{{\footnotesize\itshape Continued on next page}} \\%
\endfoot%
\hline%
\endlastfoot%
\textnormal{Steel Self-Weight Applied} & 76.4 kN/m\textsuperscript{3} \\*[6pt]
\hline
\textnormal{Concrete Deck Weight} & 24.5 kN/m\textsuperscript{3} \\*[6pt]
\hline
\textnormal{Self-Weight Factor} & 1.0 \\*[6pt]%
\end{longtable}

\vspace{1em}
\begin{longtable}{|L{5.31cm}|p{9.66cm}|}%
\caption{\textbf{Dead Load for Surfacing (DW)}}\\%
\hline%
\rowcolor{tableHeaderBg}\hspace{0pt}\textbf{parameter} & \hspace{0pt}\textbf{value} \\%
\hline%
\endfirsthead%
\caption[]{\textbf{Dead Load for Surfacing (DW)} (continued)}\\%
\hline%
\rowcolor{tableHeaderBg}\hspace{0pt}\textbf{parameter} & \hspace{0pt}\textbf{value} \\%
\hline%
\endhead%
\multicolumn{2}{r}{{\footnotesize\itshape Continued on next page}} \\%
\endfoot%
\hline%
\endlastfoot%
\textnormal{Wearing Course Load} & Concrete x 50.0 \\*[6pt]
\hline
\textnormal{Additional SIDL (Crash Barrier)} & 7.56 kN/m per barrier \\*[6pt]
\hline
\textnormal{Railing Load} & 6.806 kN/m\sdstar{} \\*[6pt]%
\end{longtable}

\vspace{1em}
\begin{longtable}{|L{4.38cm}|C{2.47cm}|C{3.99cm}|C{3.23cm}|}%
\caption{\textbf{Vehicle Live Load (LL)}}\\%
\hline%
\rowcolor{tableHeaderBg}\hspace{0pt}\textbf{Vehicle Class} & \hspace{0pt}\textbf{Total Load (kN)} & \hspace{0pt}\textbf{Impact Factor (1+IM)} & \hspace{0pt}\textbf{Braking Load Considered?} \\%
\hline%
\endfirsthead%
\caption[]{\textbf{Vehicle Live Load (LL)} (continued)}\\%
\hline%
\rowcolor{tableHeaderBg}\hspace{0pt}\textbf{Vehicle Class} & \hspace{0pt}\textbf{Total Load (kN)} & \hspace{0pt}\textbf{Impact Factor (1+IM)} & \hspace{0pt}\textbf{Braking Load Considered?} \\%
\hline%
\endhead%
\multicolumn{4}{r}{{\footnotesize\itshape Continued on next page}} \\%
\endfoot%
\hline%
\endlastfoot%
IRC Class A & 543.5 & 1.207 (IRC 6 Cl. 208.2) & As per IRC 6 Cl. 211.2 \\*[6pt]
\hline
IRC Class 70R (Wheeled) & 981.0 & 1.100 (IRC 6 Cl. 208.3) & As per IRC 6 Cl. 211.2 \\*[6pt]
\hline
IRC Class 70R (Tracked) & 686.8 & 1.100 (IRC 6 Cl. 208.3) & As per IRC 6 Cl. 211.2 \\*[6pt]
\hline
IRC Class 70R (Bogie) & ---\textsuperscript{\dag} & 1.100 (IRC 6 Cl. 208.3) & As per IRC 6 Cl. 211.2 \\*[6pt]
\hline
IRC Class AA (Wheeled) & ---\textsuperscript{\dag} & 1.100 (IRC 6 Cl. 208.3) & As per IRC 6 Cl. 211.2 \\*[6pt]
\hline
IRC Class AA (Tracked) & ---\textsuperscript{\dag} & 1.100 (IRC 6 Cl. 208.3) & As per IRC 6 Cl. 211.2 \\*[6pt]
\hline
IRC Class SV & 3776.8 & --- & Yes \\*[6pt]
\hline
IRC Fatigue Vehicle & 392.4 & --- & As per IRC 6 Cl. 211.2 \\*[6pt]
\hline
\textbf{Braking load (IRC 6 Cl. 211.2)} & \multicolumn{3}{p{10.54cm}|}{Not available --- the design pipeline's IRC 6 Cl. 211.2 braking-force function returns no value for 2 design lanes (known upstream defect D-18).} \\*[6pt]
\hline
\textbf{Braking load eccentricity (IRC 6 Cl. 211.3)} & \multicolumn{3}{p{10.54cm}|}{1.20 m above top of deck} \\*[6pt]
\hline
\textbf{Centrifugal force (IRC 6 Cl. 212)} & \multicolumn{3}{p{10.54cm}|}{Not applicable (straight bridge) - IRC 6 Cl. 212 applies to bridges on a horizontal curve; this bridge is straight in plan (no horizontal-curve input)} \\*[6pt]%
\end{longtable}
\noindent\textit{Note: Only the vehicle classes selected by the user are listed. Total load = sum of the axle / track loads of the IRC 6 Cl. 204 vehicle definition.}
\par\noindent\textit{\textsuperscript{\dag} Vehicle not defined in the OsdagBridge IRC 6 vehicle library --- total load not available.}


\vspace{1em}
\begin{longtable}{|L{5.31cm}|p{9.66cm}|}%
\caption{\textbf{Footway Live Load (FPLL)}}\\%
\hline%
\rowcolor{tableHeaderBg}\hspace{0pt}\textbf{Parameter} & \hspace{0pt}\textbf{Value} \\%
\hline%
\endfirsthead%
\caption[]{\textbf{Footway Live Load (FPLL)} (continued)}\\%
\hline%
\rowcolor{tableHeaderBg}\hspace{0pt}\textbf{Parameter} & \hspace{0pt}\textbf{Value} \\%
\hline%
\endhead%
\multicolumn{2}{r}{{\footnotesize\itshape Continued on next page}} \\%
\endfoot%
\hline%
\endlastfoot%
\textnormal{Footpath Provided} & None \\*[6pt]
\hline
\textnormal{Footway Live Load Intensity} & \textit{Not applicable (no footpath provided)} \\*[6pt]%
\end{longtable}

\vspace{1em}
\begin{longtable}{|L{5.31cm}|p{9.66cm}|}%
\caption{\textbf{Wind Load (WL) --- per IRC 6}}\\%
\hline%
\rowcolor{tableHeaderBg}\hspace{0pt}\textbf{parameter} & \hspace{0pt}\textbf{value} \\%
\hline%
\endfirsthead%
\caption[]{\textbf{Wind Load (WL) --- per IRC 6} (continued)}\\%
\hline%
\rowcolor{tableHeaderBg}\hspace{0pt}\textbf{parameter} & \hspace{0pt}\textbf{value} \\%
\hline%
\endhead%
\multicolumn{2}{r}{{\footnotesize\itshape Continued on next page}} \\%
\endfoot%
\hline%
\endlastfoot%
\textnormal{Basic Wind Speed, Vb} & 44 m/s [from Project Location] \\*[6pt]
\hline
\textnormal{Terrain Type} & Plain Terrain \\*[6pt]
\hline
\textnormal{Average Exposed Height, H (m)} & 10 m \\*[6pt]
\hline
\textnormal{Hourly Mean Wind Speed, Vz} & 37.07 m/s \\*[6pt]
\hline
\textnormal{Hourly Wind Pressure, Pz} & 824.36 N/m² \\*[6pt]
\hline
\textnormal{Transverse Wind Force} & \textit{Not available} \\*[6pt]
\hline
\textnormal{Longitudinal Wind Force} & \textit{Not available} \\*[6pt]
\hline
\textnormal{Vertical Wind Force} & \textit{Not available} \\*[6pt]%
\end{longtable}

\vspace{1em}
\begin{longtable}{|L{5.31cm}|p{9.66cm}|}%
\caption{\textbf{Earthquake Load (EL) --- per IRC 6}}\\%
\hline%
\rowcolor{tableHeaderBg}\hspace{0pt}\textbf{parameter} & \hspace{0pt}\textbf{value} \\%
\hline%
\endfirsthead%
\caption[]{\textbf{Earthquake Load (EL) --- per IRC 6} (continued)}\\%
\hline%
\rowcolor{tableHeaderBg}\hspace{0pt}\textbf{parameter} & \hspace{0pt}\textbf{value} \\%
\hline%
\endhead%
\multicolumn{2}{r}{{\footnotesize\itshape Continued on next page}} \\%
\endfoot%
\hline%
\endlastfoot%
\textnormal{Seismic Zone} & III [from Project Location] \\*[6pt]
\hline
\textnormal{Zone Factor, Z} & 0.16 \\*[6pt]
\hline
\textnormal{Importance Factor, I} & 1.0 \\*[6pt]
\hline
\textnormal{Type of Soil} & Type I – Rocky or Hard \\*[6pt]
\hline
\textnormal{Sa/g} & 2.8 \\*[6pt]
\hline
\textnormal{Horizontal Seismic Coefficient, Ah} & 0.224 \\*[6pt]
\hline
\textnormal{Vertical Seismic Coefficient, Av} & 0.1493 \\*[6pt]
\hline
\textnormal{Horizontal Seismic Force (longitudinal)} & \textit{Not available} \\*[6pt]
\hline
\textnormal{Horizontal Seismic Force (transverse)} & \textit{Not available} \\*[6pt]%
\end{longtable}

\vspace{1em}
\begin{longtable}{|L{5.31cm}|p{9.66cm}|}%
\caption{\textbf{Temperature Load (TL) --- per IRC 6}}\\%
\hline%
\rowcolor{tableHeaderBg}\hspace{0pt}\textbf{parameter} & \hspace{0pt}\textbf{value} \\%
\hline%
\endfirsthead%
\caption[]{\textbf{Temperature Load (TL) --- per IRC 6} (continued)}\\%
\hline%
\rowcolor{tableHeaderBg}\hspace{0pt}\textbf{parameter} & \hspace{0pt}\textbf{value} \\%
\hline%
\endhead%
\multicolumn{2}{r}{{\footnotesize\itshape Continued on next page}} \\%
\endfoot%
\hline%
\endlastfoot%
\textnormal{Maximum Shade Temperature} & 42.2 $^\circ$C \\*[6pt]
\hline
\textnormal{Minimum Shade Temperature} & 7.4 $^\circ$C \\*[6pt]
\hline
\textnormal{Effective Bridge Temp. Range} & -2.60 to 57.20 $^\circ$C \\*[6pt]
\hline
\textnormal{Temperature Rise / Fall for Design} & +29.90 $^\circ$C / \textminus{}29.90 $^\circ$C \\*[6pt]%
\end{longtable}

\vspace{1em}
\begin{longtable}{|C{3.86cm}|p{11.10cm}|}%
\caption{\textbf{Load Combinations}}\\%
\hline%
\rowcolor{tableHeaderBg}\hspace{0pt}\textbf{Combination ID} & \hspace{0pt}\textbf{Load Cases} \\%
\hline%
\endfirsthead%
\caption[]{\textbf{Load Combinations} (continued)}\\%
\hline%
\rowcolor{tableHeaderBg}\hspace{0pt}\textbf{Combination ID} & \hspace{0pt}\textbf{Load Cases} \\%
\hline%
\endhead%
\multicolumn{2}{r}{{\footnotesize\itshape Continued on next page}} \\%
\endfoot%
\hline%
\endlastfoot%
ULS-01 & DL(1.35/1.00) + SIDL(1.75/1.00) + LL(1.50) + WL(0.90) + TL(0.90) \\*[6pt]
\hline
ULS-02 & DL(1.35/1.00) + SIDL(1.75/1.00) + LL(1.15) + WL(1.50) + TL(0.90) \\*[6pt]
\hline
ULS-03 & DL(1.35/1.00) + SIDL(1.75/1.00) + LL(1.15) + WL(0.90) + TL(1.50) \\*[6pt]
\hline
ULS-04 & DL(1.00/1.00) + SIDL(1.00/1.00) + LL(0.75) + TL(0.50) + VC(1.00) \\*[6pt]
\hline
ULS-05 & DL(1.00/1.00) + SIDL(1.00/1.00) + LL(0.75) + TL(0.50) + BI(1.00) \\*[6pt]
\hline
ULS-06 & DL(1.00/1.00) + SIDL(1.00/1.00) + LL(0.75) + TL(0.50) + FB(1.00) \\*[6pt]
\hline
ULS-07 & DL(1.35/1.00) + SIDL(1.75/1.00) + LL(0.20) + TL(0.50) + EQ(1.50) \\*[6pt]
\hline
ULS-08 & DL(1.35/1.00) + SIDL(1.75/1.00) + LL(0.20) + TL(0.50) + EQ(0.75) \\*[6pt]
\hline
SLS-01 & DL(1.00) + SIDL(1.20/1.00) + LL(1.00) + WL(0.60) + TL(0.60) \\*[6pt]
\hline
SLS-02 & DL(1.00) + SIDL(1.20/1.00) + LL(0.75) + WL(1.00) + TL(0.60) \\*[6pt]
\hline
SLS-03 & DL(1.00) + SIDL(1.20/1.00) + LL(0.75) + WL(0.60) + TL(1.00) \\*[6pt]
\hline
SLS-04 & DL(1.00) + SIDL(1.20/1.00) + LL(0.75) + WL(0.50) + TL(0.50) \\*[6pt]
\hline
SLS-05 & DL(1.00) + SIDL(1.20/1.00) + LL(0.20) + WL(0.60) + TL(0.50) \\*[6pt]
\hline
SLS-06 & DL(1.00) + SIDL(1.20/1.00) + LL(0.20) + WL(0.50) + TL(0.60) \\*[6pt]
\hline
SLS-07 & DL(1.00) + SIDL(1.20/1.00) + LL(0.00) + WL(0.00) + TL(0.50) \\*[6pt]%
\end{longtable}

\noindent\textit{Note: All IRC 6 load combinations are auto-generated by OsdagBridge. User-defined custom combinations, if any, are appended.}


\chapter{Analysis Results}

A grillage model was used for structural analysis. The deck is idealized as a grid of elastic beam elements --- longitudinal members represent the composite steel girders with effective slab, and transverse members represent the slab or cross frames. This section summarizes the critical output from that analysis.

\vspace{1em}
\begingroup
\footnotesize
\setlength{\tabcolsep}{3pt}
\renewcommand{\arraystretch}{1.25}

\begin{longtable}{|>{\centering\arraybackslash}p{2.90cm}|>{\centering\arraybackslash}p{1.59cm}|>{\centering\arraybackslash}p{1.12cm}|>{\centering\arraybackslash}p{1.22cm}|>{\centering\arraybackslash}p{1.59cm}|>{\centering\arraybackslash}p{1.12cm}|>{\centering\arraybackslash}p{1.22cm}|>{\centering\arraybackslash}p{1.50cm}|>{\centering\arraybackslash}p{1.50cm}|}%
\caption{\textbf{Summary of Maximum Demands}}\\%
\hline%
\multirow{2}{*}{\makecell{\hspace{0pt}\textbf{Load}\\\hspace{0pt}\textbf{Case/}\\\hspace{0pt}\textbf{Comb.}}}
& \multicolumn{3}{c|}{\hspace{0pt}\textbf{Bending Moment}}
& \multicolumn{3}{c|}{\hspace{0pt}\textbf{Shear Force}}
& \multicolumn{2}{c|}{\makecell{\hspace{0pt}\textbf{Reaction at}\\\hspace{0pt}\textbf{Supports}}}\\
\cline{2-9}
& \makecell{\hspace{0pt}\textbf{Max}\\\hspace{0pt}\textbf{(kNm)}}
& \hspace{0pt}\textbf{Girder}
& \makecell{\hspace{0pt}\textbf{Loc.}\\\hspace{0pt}\textbf{(m)}}
& \makecell{\hspace{0pt}\textbf{Max}\\\hspace{0pt}\textbf{(kN)}}
& \hspace{0pt}\textbf{Girder}
& \makecell{\hspace{0pt}\textbf{Loc.}\\\hspace{0pt}\textbf{(m)}}
& \makecell{\hspace{0pt}\textbf{Left}}
& \makecell{\hspace{0pt}\textbf{Right}} \\%
\hline%
\endfirsthead%
\caption[]{\textbf{Summary of Maximum Demands} (continued)}\\%
\hline%
\multirow{2}{*}{\makecell{\hspace{0pt}\textbf{Load}\\\hspace{0pt}\textbf{Case/}\\\hspace{0pt}\textbf{Comb.}}}
& \multicolumn{3}{c|}{\hspace{0pt}\textbf{Bending Moment}}
& \multicolumn{3}{c|}{\hspace{0pt}\textbf{Shear Force}}
& \multicolumn{2}{c|}{\makecell{\hspace{0pt}\textbf{Reaction at}\\\hspace{0pt}\textbf{Supports}}}\\
\cline{2-9}
& \makecell{\hspace{0pt}\textbf{Max}\\\hspace{0pt}\textbf{(kNm)}}
& \hspace{0pt}\textbf{Girder}
& \makecell{\hspace{0pt}\textbf{Loc.}\\\hspace{0pt}\textbf{(m)}}
& \makecell{\hspace{0pt}\textbf{Max}\\\hspace{0pt}\textbf{(kN)}}
& \hspace{0pt}\textbf{Girder}
& \makecell{\hspace{0pt}\textbf{Loc.}\\\hspace{0pt}\textbf{(m)}}
& \makecell{\hspace{0pt}\textbf{Left}}
& \makecell{\hspace{0pt}\textbf{Right}} \\%
\hline%
\endhead%
\multicolumn{9}{r}{{\footnotesize\itshape Continued on next page}} \\%
\endfoot%
\hline%
\endlastfoot%
SLS\_\allowbreak{}RARE\_\allowbreak{}3:\allowbreak{} 1.0DL + 1.2DW + 0.75LL + 0.6WL & 2115.421 & G2 & 13.500 & 361.509 & G1 & 0.000 & -1244.723 & 977.910 \\*[6pt]
\hline
ACCIDENTAL\_\allowbreak{}2:\allowbreak{} 1.0DL + 1.0DW + 0.75LL & 2224.950 & G2 & 13.500 & 379.980 & G1 & 0.000 & -1308.880 & 1042.068 \\*[6pt]
\hline
ACCIDENTAL\_\allowbreak{}3:\allowbreak{} 1.0DL + 1.0DW + 0.75LL & 2224.950 & G2 & 13.500 & 379.980 & G1 & 0.000 & -1308.880 & 1042.068 \\*[6pt]
\hline
BASIC\_\allowbreak{}1:\allowbreak{} 1.35DL + 1.75DW + 1.5LL + 0.9WL & 2115.421 & G2 & 13.500 & 361.509 & G1 & 0.000 & -1244.723 & 977.910 \\*[6pt]
\hline
1.5 EQ (a) & 94.042 & G1 & 15.000 & 15.683 & G1 & 0.000 & 54.473 & -54.473 \\*[6pt]
\hline
SLS\_\allowbreak{}FREQUENT\_\allowbreak{}2:\allowbreak{} 1.0DL + 1.2DW + 0.2LL + 0.6WL & 2115.421 & G2 & 13.500 & 361.509 & G1 & 0.000 & -1244.723 & 977.910 \\*[6pt]
\hline
BASIC\_\allowbreak{}3:\allowbreak{} 1.35DL + 1.75DW + 1.15LL + 0.9WL & 2115.421 & G2 & 13.500 & 361.509 & G1 & 0.000 & -1244.723 & 977.910 \\*[6pt]
\hline
SLS\_\allowbreak{}QP\_\allowbreak{}2:\allowbreak{} 1.0DL + 1.0DW & 1402.699 & G1 & 15.000 & 237.286 & G1 & 0.000 & -820.260 & 820.260 \\*[6pt]
\hline
SLS\_\allowbreak{}FREQUENT\_\allowbreak{}4:\allowbreak{} 1.0DL + 1.0DW + 0.75LL + 0.5WL & 2115.421 & G2 & 13.500 & 361.509 & G1 & 0.000 & -1244.723 & 977.910 \\*[6pt]
\hline
ACCIDENTAL\_\allowbreak{}1:\allowbreak{} 1.0DL + 1.0DW + 0.75LL & 2224.950 & G2 & 13.500 & 379.980 & G1 & 0.000 & -1308.880 & 1042.068 \\*[6pt]
\hline
BASIC\_\allowbreak{}5:\allowbreak{} 1.0DL + 1.0DW + 1.15LL + 1.5WL & 2115.421 & G2 & 13.500 & 361.509 & G1 & 0.000 & -1244.723 & 977.910 \\*[6pt]
\hline
SLS\_\allowbreak{}QP\_\allowbreak{}1:\allowbreak{} 1.0DL + 1.2DW & 1402.699 & G1 & 15.000 & 237.286 & G1 & 0.000 & -820.260 & 820.260 \\*[6pt]
\hline
Envelope ULS & 2224.950 & G2 & 13.500 & 379.980 & G1 & 0.000 & -1308.880 & 1042.068 \\*[6pt]
\hline
SEISMIC\_\allowbreak{}3:\allowbreak{} 1.0DL + 1.0DW + 0.2LL + 1.5EL & 2018.294 & G2 & 13.500 & 345.128 & G1 & 0.000 & -1187.830 & 921.017 \\*[6pt]
\hline
Envelope SLS & 2115.421 & G2 & 13.500 & 361.509 & G1 & 0.000 & -1244.723 & 977.910 \\*[6pt]
\hline
SEISMIC\_\allowbreak{}2:\allowbreak{} 1.35DL + 1.75DW + 0.2LL + 0.75EL & 2018.294 & G2 & 13.500 & 345.128 & G1 & 0.000 & -1187.830 & 921.017 \\*[6pt]
\hline
SLS\_\allowbreak{}FREQUENT\_\allowbreak{}3:\allowbreak{} 1.0DL + 1.2DW + 0.2LL + 0.5WL & 2115.421 & G2 & 13.500 & 361.509 & G1 & 0.000 & -1244.723 & 977.910 \\*[6pt]
\hline
1.0 DW & 199.547 & G1 & 15.000 & 33.869 & G1 & 0.000 & -116.853 & 116.853 \\*[6pt]
\hline
EQ\_\allowbreak{}Z & 0.000 & — & --- & 0.000 & — & --- & 0.000 & -0.000 \\*[6pt]
\hline
1.0 WL & 110.762 & G1 & 15.000 & 18.472 & G1 & 28.500 & 64.158 & -64.158 \\*[6pt]
\hline
SLS\_\allowbreak{}RARE\_\allowbreak{}6:\allowbreak{} 1.0DL + 1.0DW + 0.75LL + 0.6WL & 2115.421 & G2 & 13.500 & 361.509 & G1 & 0.000 & -1244.723 & 977.910 \\*[6pt]
\hline
SLS\_\allowbreak{}RARE\_\allowbreak{}1:\allowbreak{} 1.0DL + 1.2DW + 1.0LL + 0.6WL & 2115.421 & G2 & 13.500 & 361.509 & G1 & 0.000 & -1244.723 & 977.910 \\*[6pt]
\hline
1.5 EQ (b) & 94.042 & G1 & 15.000 & 15.683 & G1 & 0.000 & 54.473 & -54.473 \\*[6pt]
\hline
SLS\_\allowbreak{}FREQUENT\_\allowbreak{}1:\allowbreak{} 1.0DL + 1.2DW + 0.75LL + 0.5WL & 2115.421 & G2 & 13.500 & 361.509 & G1 & 0.000 & -1244.723 & 977.910 \\*[6pt]
\hline
SW & 205.417 & G1 & 15.000 & 34.072 & G1 & 28.500 & -119.141 & 119.141 \\*[6pt]
\hline
EQ\_\allowbreak{}Y & 208.982 & G1 & 15.000 & 34.852 & G1 & 28.500 & 121.050 & -121.050 \\*[6pt]
\hline
1.5 EQ (c) & 313.472 & G1 & 15.000 & 52.278 & G1 & 28.500 & 181.575 & -181.575 \\*[6pt]
\hline
SLS\_\allowbreak{}FREQUENT\_\allowbreak{}5:\allowbreak{} 1.0DL + 1.0DW + 0.2LL + 0.6WL & 2115.421 & G2 & 13.500 & 361.509 & G1 & 0.000 & -1244.723 & 977.910 \\*[6pt]
\hline
Crash barrier load & 162.535 & G2 & 15.000 & 28.678 & G1 & 28.500 & -87.634 & 92.345 \\*[6pt]
\hline
DD & 997.735 & G1 & 15.000 & 169.345 & G1 & 28.500 & -584.266 & 584.266 \\*[6pt]
\hline
EQ\_\allowbreak{}X & 0.000 & — & --- & 0.000 & — & --- & 0.000 & 0.000 \\*[6pt]
\hline
SEISMIC\_\allowbreak{}4:\allowbreak{} 1.0DL + 1.0DW + 0.2LL + 0.75EL & 2018.294 & G2 & 13.500 & 345.128 & G1 & 0.000 & -1187.830 & 921.017 \\*[6pt]
\hline
WL Transverse & 0.000 & — & --- & 0.000 & — & --- & -0.000 & -0.000 \\*[6pt]
\hline
SLS\_\allowbreak{}RARE\_\allowbreak{}4:\allowbreak{} 1.0DL + 1.0DW + 1.0LL + 0.6WL & 2115.421 & G2 & 13.500 & 361.509 & G1 & 0.000 & -1244.723 & 977.910 \\*[6pt]
\hline
SLS\_\allowbreak{}RARE\_\allowbreak{}5:\allowbreak{} 1.0DL + 1.0DW + 0.75LL + 1.0WL & 2115.421 & G2 & 13.500 & 361.509 & G1 & 0.000 & -1244.723 & 977.910 \\*[6pt]
\hline
WL Longitudinal & 0.000 & — & --- & 0.000 & — & --- & -0.000 & 0.000 \\*[6pt]
\hline
WL Uplift & 110.762 & G1 & 15.000 & 18.472 & G1 & 28.500 & 64.158 & -64.158 \\*[6pt]
\hline
BASIC\_\allowbreak{}4:\allowbreak{} 1.0DL + 1.0DW + 1.5LL + 0.9WL & 2115.421 & G2 & 13.500 & 361.509 & G1 & 0.000 & -1244.723 & 977.910 \\*[6pt]
\hline
1.0 DL & 1203.152 & G1 & 15.000 & 203.417 & G1 & 28.500 & -703.407 & 703.407 \\*[6pt]
\hline
BASIC\_\allowbreak{}6:\allowbreak{} 1.0DL + 1.0DW + 1.15LL + 0.9WL & 2115.421 & G2 & 13.500 & 361.509 & G1 & 0.000 & -1244.723 & 977.910 \\*[6pt]
\hline
1.0 DL + 1.0 LL & 2027.623 & G2 & 13.500 & 346.111 & G1 & 0.000 & -1192.027 & 925.214 \\*[6pt]
\hline
SLS\_\allowbreak{}FREQUENT\_\allowbreak{}6:\allowbreak{} 1.0DL + 1.0DW + 0.2LL + 0.5WL & 2115.421 & G2 & 13.500 & 361.509 & G1 & 0.000 & -1244.723 & 977.910 \\*[6pt]
\hline
BASIC\_\allowbreak{}2:\allowbreak{} 1.35DL + 1.75DW + 1.15LL + 1.5WL & 2115.421 & G2 & 13.500 & 361.509 & G1 & 0.000 & -1244.723 & 977.910 \\*[6pt]
\hline
SLS\_\allowbreak{}RARE\_\allowbreak{}2:\allowbreak{} 1.0DL + 1.2DW + 0.75LL + 1.0WL & 2115.421 & G2 & 13.500 & 361.509 & G1 & 0.000 & -1244.723 & 977.910 \\*[6pt]
\hline
SEISMIC\_\allowbreak{}1:\allowbreak{} 1.35DL + 1.75DW + 0.2LL + 1.5EL & 2018.294 & G2 & 13.500 & 345.128 & G1 & 0.000 & -1187.830 & 921.017 \\*[6pt]%
\end{longtable}
\endgroup
\noindent\textit{Note: The girder design results contain 6 longitudinal members (G1-G6) although the bridge has 4 girders. The grillage model adds a longitudinal line at each deck edge, and the design module designs and labels these edge lines as girders as well. Results are reported with the labels used by the design module, so these labels do not necessarily denote the girders numbered in Chapter 2. This is an issue of the analysis/design module (not changed by the report) and requires engineering verification.}


\vspace{1em}
\begin{longtable}{|>{\centering\arraybackslash}p{4.84cm}|>{\centering\arraybackslash}p{4.84cm}|>{\centering\arraybackslash}p{4.84cm}|}%
\caption{\textbf{Reactions at Supports --- Governing Values}}\\%
\hline%
\rowcolor{tableHeaderBg}\hspace{0pt}\textbf{Support} & \hspace{0pt}\textbf{Governing Load Case} & \hspace{0pt}\textbf{Governing Reaction (kN), max.\ $|R|$} \\%
\hline%
\endfirsthead%
\caption[]{\textbf{Reactions at Supports --- Governing Values} (continued)}\\%
\hline%
\rowcolor{tableHeaderBg}\hspace{0pt}\textbf{Support} & \hspace{0pt}\textbf{Governing Load Case} & \hspace{0pt}\textbf{Governing Reaction (kN), max.\ $|R|$} \\%
\hline%
\endhead%
\multicolumn{3}{r}{{\footnotesize\itshape Continued on next page}} \\%
\endfoot%
\hline%
\endlastfoot%
Left support & ACCIDENTAL\_\allowbreak{}1:\allowbreak{} 1.0DL + 1.0DW + 0.75LL \textit{(and 3 other load cases with the same value)} & -1308.880 \\*[6pt]
\hline
Right support & ACCIDENTAL\_\allowbreak{}1:\allowbreak{} 1.0DL + 1.0DW + 0.75LL \textit{(and 3 other load cases with the same value)} & 1042.068 \\*[6pt]%
\end{longtable}

\vspace{1em}
\begin{longtable}{|L{6.76cm}|p{8.21cm}|}%
\caption{\textbf{Deflection Summary (Live Load \& Total Load)}}\\%
\hline%
\rowcolor{tableHeaderBg}\hspace{0pt}\textbf{parameter} & \hspace{0pt}\textbf{value} \\%
\hline%
\endfirsthead%
\caption[]{\textbf{Deflection Summary (Live Load \& Total Load)} (continued)}\\%
\hline%
\rowcolor{tableHeaderBg}\hspace{0pt}\textbf{parameter} & \hspace{0pt}\textbf{value} \\%
\hline%
\endhead%
\multicolumn{2}{r}{{\footnotesize\itshape Continued on next page}} \\%
\endfoot%
\hline%
\endlastfoot%
\textnormal{Deflection due to Live Load, $\delta_{LL}$} & 19.385 mm \\*[6pt]
\hline
\textnormal{Allowable Live Load Deflection ($\Delta_{allow}$)} & L/800 = 37.5 mm \\*[6pt]
\hline
\textnormal{Live Load Deflection Check Status} & PASS \\*[6pt]
\hline
\textnormal{Deflection due to Total Load, $\delta_{total}$} & 51.225 mm \\*[6pt]
\hline
\textnormal{Allowable Total Deflection ($\Delta_{allow}$)} & L/600 = 50.0 mm \\*[6pt]
\hline
\textnormal{Total Load Deflection Check Status} & \textcolor{red}{FAIL} \\*[6pt]%
\end{longtable}
\noindent\textit{Note: Deflection:\allowbreak{} the per-girder load-effect results used in Table 4.3 give 19.385 mm (live) and 51.225 mm (total), whereas the IRC 22 deflection checks of the girder design module (Section 5.1) use 9.631 mm and 22.871 mm. The two calculations of the application disagree, so the deflection status differs between Table 4.3 and Section 5. Requires engineering verification.}


\vspace{1em}
\noindent
\begin{figure}[H]
\vspace{-0.5em}
\centering
\includegraphics[width=0.75\textwidth]{C:/Users/udayk/AppData/Local/Temp/tmpymqxq2pd/images/bm_envelope.png}
\vspace{0.5em}
\caption*{\small Bending Moment Envelope (Envelope ULS): Max/min BM along span. X-axis: distance from left support (m). Y-axis: Bending Moment (kN-m).}
\vspace{-0.5em}
\end{figure}
\begin{figure}[H]
\vspace{-0.5em}
\centering
\includegraphics[width=0.75\textwidth]{C:/Users/udayk/AppData/Local/Temp/tmpymqxq2pd/images/sf_envelope.png}
\vspace{0.5em}
\caption*{\small Shear Force Envelope (Envelope ULS): Max/min SF along span. X-axis: distance from left support (m). Y-axis: Shear Force (kN).}
\vspace{-0.5em}
\end{figure}
\begin{figure}[H]
\vspace{-0.5em}
\centering
\includegraphics[width=0.75\textwidth]{C:/Users/udayk/AppData/Local/Temp/tmpymqxq2pd/images/defl_ll.png}
\vspace{0.5em}
\caption*{\small Vertical Deflection D$_y$ (1.0 LL): Maximum deflection along span. Load Case: 1.0 LL, Combination: $D_y$. Nodes shown. Isometric view.}
\vspace{-0.5em}
\end{figure}


\chapter{Design Checks}

This section presents all structural design checks performed by OsdagBridge. For each member, the demand from the governing load combination, the code-based capacity, and the utilization ratio are tabulated. All checks reference IS 800:2007 and IRC 22:2014 unless stated otherwise.

\par\medskip\noindent\fbox{\parbox{0.97\textwidth}{\small\textbf{Data consistency notes:}\begin{itemize}\setlength{\itemsep}{1pt}
\item The girder design results contain 6 longitudinal members (G1-G6) although the bridge has 4 girders. The grillage model adds a longitudinal line at each deck edge, and the design module designs and labels these edge lines as girders as well. Results are reported with the labels used by the design module, so these labels do not necessarily denote the girders numbered in Chapter 2. This is an issue of the analysis/design module (not changed by the report) and requires engineering verification.
\item Girder section properties:\allowbreak{} the grillage analysis model uses the stored section area 349.2 cm2, whereas the final designed section (1670 x 510 x 22 x 510 x 22) has 387.0 cm2. Member forces were therefore computed with properties that differ from the checked section. Requires engineering verification.
\item Cross-bracing spacing:\allowbreak{} the input lists 15 m (Table 2.9), but the girder design checks and the cross-bracing force calculation used 3 m. The design modules read a spacing value that is not set by the input, so their default was applied. Requires engineering verification.
\item Deflection:\allowbreak{} the per-girder load-effect results used in Table 4.3 give 19.385 mm (live) and 51.225 mm (total), whereas the IRC 22 deflection checks of the girder design module (Section 5.1) use 9.631 mm and 22.871 mm. The two calculations of the application disagree, so the deflection status differs between Table 4.3 and Section 5. Requires engineering verification.
\end{itemize}}}\par\medskip


% Table 5.1 starts with an unbreakable 13-row girder group, so the section
% starts on a new page (its heading was left alone at a page bottom).
% \clearpage, not \needspace: a \needspace break before a longtable is
% processed by longtable's output routine and misplaces its headings.
\clearpage
\section{Plate Girder Design}
\label{sec:plate-girder}

\vspace{1em}
\begin{longtable}{|C{2.34cm}|L{7.50cm}|>{\centering\arraybackslash}p{4.68cm}|}%
\caption{\textbf{Girder Section Properties (Final Optimized / User-selected)}}\\%
\hline%
\rowcolor{tableHeaderBg}\hspace{0pt}\textbf{Girder} & \hspace{0pt}\textbf{Property} & \hspace{0pt}\textbf{Value} \\%
\hline%
\endfirsthead%
\caption[]{\textbf{Girder Section Properties (Final Optimized / User-selected)} (continued)}\\%
\hline%
\rowcolor{tableHeaderBg}\hspace{0pt}\textbf{Girder} & \hspace{0pt}\textbf{Property} & \hspace{0pt}\textbf{Value} \\%
\hline%
\endhead%
\multicolumn{3}{r}{{\footnotesize\itshape Continued on next page}} \\%
\endfoot%
\hline%
\endlastfoot%
\multirow{13}{*}{\makecell{G1}} & \textnormal{Depth, D (mm)} & 1670.0 \\*[6pt]
\cline{2-3}\noalign{\nobreak}
 & \textnormal{Top Flange Width, $b_f$ (mm)} & 510.0 \\*[6pt]
\cline{2-3}\noalign{\nobreak}
 & \textnormal{Bottom Flange Width, $b_f$ (mm)} & 510.0 \\*[6pt]
\cline{2-3}\noalign{\nobreak}
 & \textnormal{Top Flange Thickness, $t_f$ (mm)} & 22.0 \\*[6pt]
\cline{2-3}\noalign{\nobreak}
 & \textnormal{Bottom Flange Thickness, $t_f$ (mm)} & 22.0 \\*[6pt]
\cline{2-3}\noalign{\nobreak}
 & \textnormal{Web Thickness, $t_w$ (mm)} & 10.0 \\*[6pt]
\cline{2-3}\noalign{\nobreak}
 & \textnormal{Gross Area, A (cm$^2$)} & 387.0 \\*[6pt]
\cline{2-3}\noalign{\nobreak}
 & \textnormal{Moment of Inertia, $I_z$ (cm$^4$)} & 1881957.85 \\*[6pt]
\cline{2-3}\noalign{\nobreak}
 & \textnormal{Elastic Section Modulus, $Z_{ez}$ (cm$^3$)} & 22538.42 \\*[6pt]
\cline{2-3}\noalign{\nobreak}
 & \textnormal{Plastic Section Modulus, $Z_{pz}$ (cm$^3$)} & 25100.25 \\*[6pt]
\cline{2-3}\noalign{\nobreak}
 & \textnormal{Effective Slab Width, $b_{eff}$ (mm)} & 2100.0 \\*[6pt]
\cline{2-3}\noalign{\nobreak}
 & \textnormal{Transformed Composite $I_z$ (cm$^4$)} & 4215209.56 \\*[6pt]
\cline{2-3}\noalign{\nobreak}
 & \textnormal{Depth to Plastic Neutral Axis (mm)} & 264.8 \\*[6pt]
\hline
\multirow{13}{*}{\makecell{G2}} & \textnormal{Depth, D (mm)} & 1670.0 \\*[6pt]
\cline{2-3}\noalign{\nobreak}
 & \textnormal{Top Flange Width, $b_f$ (mm)} & 510.0 \\*[6pt]
\cline{2-3}\noalign{\nobreak}
 & \textnormal{Bottom Flange Width, $b_f$ (mm)} & 510.0 \\*[6pt]
\cline{2-3}\noalign{\nobreak}
 & \textnormal{Top Flange Thickness, $t_f$ (mm)} & 22.0 \\*[6pt]
\cline{2-3}\noalign{\nobreak}
 & \textnormal{Bottom Flange Thickness, $t_f$ (mm)} & 22.0 \\*[6pt]
\cline{2-3}\noalign{\nobreak}
 & \textnormal{Web Thickness, $t_w$ (mm)} & 10.0 \\*[6pt]
\cline{2-3}\noalign{\nobreak}
 & \textnormal{Gross Area, A (cm$^2$)} & 387.0 \\*[6pt]
\cline{2-3}\noalign{\nobreak}
 & \textnormal{Moment of Inertia, $I_z$ (cm$^4$)} & 1881957.85 \\*[6pt]
\cline{2-3}\noalign{\nobreak}
 & \textnormal{Elastic Section Modulus, $Z_{ez}$ (cm$^3$)} & 22538.42 \\*[6pt]
\cline{2-3}\noalign{\nobreak}
 & \textnormal{Plastic Section Modulus, $Z_{pz}$ (cm$^3$)} & 25100.25 \\*[6pt]
\cline{2-3}\noalign{\nobreak}
 & \textnormal{Effective Slab Width, $b_{eff}$ (mm)} & 2100.0 \\*[6pt]
\cline{2-3}\noalign{\nobreak}
 & \textnormal{Transformed Composite $I_z$ (cm$^4$)} & 4215209.56 \\*[6pt]
\cline{2-3}\noalign{\nobreak}
 & \textnormal{Depth to Plastic Neutral Axis (mm)} & 264.8 \\*[6pt]
\hline
\multirow{13}{*}{\makecell{G3}} & \textnormal{Depth, D (mm)} & 1670.0 \\*[6pt]
\cline{2-3}\noalign{\nobreak}
 & \textnormal{Top Flange Width, $b_f$ (mm)} & 510.0 \\*[6pt]
\cline{2-3}\noalign{\nobreak}
 & \textnormal{Bottom Flange Width, $b_f$ (mm)} & 510.0 \\*[6pt]
\cline{2-3}\noalign{\nobreak}
 & \textnormal{Top Flange Thickness, $t_f$ (mm)} & 22.0 \\*[6pt]
\cline{2-3}\noalign{\nobreak}
 & \textnormal{Bottom Flange Thickness, $t_f$ (mm)} & 22.0 \\*[6pt]
\cline{2-3}\noalign{\nobreak}
 & \textnormal{Web Thickness, $t_w$ (mm)} & 10.0 \\*[6pt]
\cline{2-3}\noalign{\nobreak}
 & \textnormal{Gross Area, A (cm$^2$)} & 387.0 \\*[6pt]
\cline{2-3}\noalign{\nobreak}
 & \textnormal{Moment of Inertia, $I_z$ (cm$^4$)} & 1881957.85 \\*[6pt]
\cline{2-3}\noalign{\nobreak}
 & \textnormal{Elastic Section Modulus, $Z_{ez}$ (cm$^3$)} & 22538.42 \\*[6pt]
\cline{2-3}\noalign{\nobreak}
 & \textnormal{Plastic Section Modulus, $Z_{pz}$ (cm$^3$)} & 25100.25 \\*[6pt]
\cline{2-3}\noalign{\nobreak}
 & \textnormal{Effective Slab Width, $b_{eff}$ (mm)} & 2100.0 \\*[6pt]
\cline{2-3}\noalign{\nobreak}
 & \textnormal{Transformed Composite $I_z$ (cm$^4$)} & 4215209.56 \\*[6pt]
\cline{2-3}\noalign{\nobreak}
 & \textnormal{Depth to Plastic Neutral Axis (mm)} & 264.8 \\*[6pt]
\hline
\multirow{13}{*}{\makecell{G4}} & \textnormal{Depth, D (mm)} & 1670.0 \\*[6pt]
\cline{2-3}\noalign{\nobreak}
 & \textnormal{Top Flange Width, $b_f$ (mm)} & 510.0 \\*[6pt]
\cline{2-3}\noalign{\nobreak}
 & \textnormal{Bottom Flange Width, $b_f$ (mm)} & 510.0 \\*[6pt]
\cline{2-3}\noalign{\nobreak}
 & \textnormal{Top Flange Thickness, $t_f$ (mm)} & 22.0 \\*[6pt]
\cline{2-3}\noalign{\nobreak}
 & \textnormal{Bottom Flange Thickness, $t_f$ (mm)} & 22.0 \\*[6pt]
\cline{2-3}\noalign{\nobreak}
 & \textnormal{Web Thickness, $t_w$ (mm)} & 10.0 \\*[6pt]
\cline{2-3}\noalign{\nobreak}
 & \textnormal{Gross Area, A (cm$^2$)} & 387.0 \\*[6pt]
\cline{2-3}\noalign{\nobreak}
 & \textnormal{Moment of Inertia, $I_z$ (cm$^4$)} & 1881957.85 \\*[6pt]
\cline{2-3}\noalign{\nobreak}
 & \textnormal{Elastic Section Modulus, $Z_{ez}$ (cm$^3$)} & 22538.42 \\*[6pt]
\cline{2-3}\noalign{\nobreak}
 & \textnormal{Plastic Section Modulus, $Z_{pz}$ (cm$^3$)} & 25100.25 \\*[6pt]
\cline{2-3}\noalign{\nobreak}
 & \textnormal{Effective Slab Width, $b_{eff}$ (mm)} & 2100.0 \\*[6pt]
\cline{2-3}\noalign{\nobreak}
 & \textnormal{Transformed Composite $I_z$ (cm$^4$)} & 4215209.56 \\*[6pt]
\cline{2-3}\noalign{\nobreak}
 & \textnormal{Depth to Plastic Neutral Axis (mm)} & 264.8 \\*[6pt]%
\end{longtable}

\vspace{1em}
\begin{longtable}{|C{2.20cm}|L{2.64cm}|C{3.08cm}|C{2.20cm}|>{\centering\arraybackslash}p{3.52cm}|}%
\caption{\textbf{Girder Section Classification}}\\%
\hline%
\rowcolor{tableHeaderBg}\hspace{0pt}\textbf{} & \hspace{0pt}\textbf{Element} & \hspace{0pt}\textbf{Slenderness Ratio} & \hspace{0pt}\textbf{Class Limit} & \hspace{0pt}\textbf{Classification} \\%
\hline%
\endfirsthead%
\caption[]{\textbf{Girder Section Classification} (continued)}\\%
\hline%
\rowcolor{tableHeaderBg}\hspace{0pt}\textbf{} & \hspace{0pt}\textbf{Element} & \hspace{0pt}\textbf{Slenderness Ratio} & \hspace{0pt}\textbf{Class Limit} & \hspace{0pt}\textbf{Classification} \\%
\hline%
\endhead%
\multicolumn{5}{r}{{\footnotesize\itshape Continued on next page}} \\%
\endfoot%
\hline%
\endlastfoot%
\multirow{3}{*}{\makecell{G1}} & Top Flange & 11.36 & 11.49 & Semi-Compact \\*[6pt]
\cline{2-5}\noalign{\nobreak}
 & Web & 162.6 & 106.49 & Slender \\*[6pt]
\cline{2-5}\noalign{\nobreak}
 & Overall Section & --- & --- & Slender \\*[6pt]
\hline
\multirow{3}{*}{\makecell{G2}} & Top Flange & 11.36 & 11.49 & Semi-Compact \\*[6pt]
\cline{2-5}\noalign{\nobreak}
 & Web & 162.6 & 106.49 & Slender \\*[6pt]
\cline{2-5}\noalign{\nobreak}
 & Overall Section & --- & --- & Slender \\*[6pt]
\hline
\multirow{3}{*}{\makecell{G3}} & Top Flange & 11.36 & 11.49 & Semi-Compact \\*[6pt]
\cline{2-5}\noalign{\nobreak}
 & Web & 162.6 & 106.49 & Slender \\*[6pt]
\cline{2-5}\noalign{\nobreak}
 & Overall Section & --- & --- & Slender \\*[6pt]
\hline
\multirow{3}{*}{\makecell{G4}} & Top Flange & 11.36 & 11.49 & Semi-Compact \\*[6pt]
\cline{2-5}\noalign{\nobreak}
 & Web & 162.6 & 106.49 & Slender \\*[6pt]
\cline{2-5}\noalign{\nobreak}
 & Overall Section & --- & --- & Slender \\*[6pt]%
\end{longtable}
\noindent\textit{Note: IS 800:2007 Table 2}

\vspace{1em}
\begin{longtable}{|C{2.20cm}|C{3.08cm}|C{3.08cm}|>{\centering\arraybackslash}p{3.70cm}|C{1.58cm}|}%
\caption{\textbf{Moment Capacity Check}}\\%
\hline%
\rowcolor{tableHeaderBg}\hspace{0pt}\textbf{} & \hspace{0pt}\textbf{Parameter} & \hspace{0pt}\textbf{Formula} & \hspace{0pt}\textbf{Value} & \hspace{0pt}\textbf{Status} \\%
\hline%
\endfirsthead%
\caption[]{\textbf{Moment Capacity Check} (continued)}\\%
\hline%
\rowcolor{tableHeaderBg}\hspace{0pt}\textbf{} & \hspace{0pt}\textbf{Parameter} & \hspace{0pt}\textbf{Formula} & \hspace{0pt}\textbf{Value} & \hspace{0pt}\textbf{Status} \\%
\hline%
\endhead%
\multicolumn{5}{r}{{\footnotesize\itshape Continued on next page}} \\%
\endfoot%
\hline%
\endlastfoot%
\multirow{3}{*}{\makecell{G1}} & Applied Moment, $M_u$ & Governing LC (ULS) & 2224.95 kN-m & --- \\*[6pt]
\cline{2-5}\noalign{\nobreak}
 & Design Moment Capacity, $M_d$ & IRC 22 Cl. 603.3.1 & 11371.85 kN-m & --- \\*[6pt]
\cline{2-5}\noalign{\nobreak}
 & Utilization Ratio, $M_u / M_d$ & $M_u / M_d$ & 0.20 & PASS \\*[6pt]
\hline
\multirow{3}{*}{\makecell{G2}} & Applied Moment, $M_u$ & Governing LC (ULS) & 2224.95 kN-m & --- \\*[6pt]
\cline{2-5}\noalign{\nobreak}
 & Design Moment Capacity, $M_d$ & IRC 22 Cl. 603.3.1 & 11371.85 kN-m & --- \\*[6pt]
\cline{2-5}\noalign{\nobreak}
 & Utilization Ratio, $M_u / M_d$ & $M_u / M_d$ & 0.20 & PASS \\*[6pt]
\hline
\multirow{3}{*}{\makecell{G3}} & Applied Moment, $M_u$ & Governing LC (ULS) & 2224.95 kN-m & --- \\*[6pt]
\cline{2-5}\noalign{\nobreak}
 & Design Moment Capacity, $M_d$ & IRC 22 Cl. 603.3.1 & 11371.85 kN-m & --- \\*[6pt]
\cline{2-5}\noalign{\nobreak}
 & Utilization Ratio, $M_u / M_d$ & $M_u / M_d$ & 0.20 & PASS \\*[6pt]
\hline
\multirow{3}{*}{\makecell{G4}} & Applied Moment, $M_u$ & Governing LC (ULS) & 2224.95 kN-m & --- \\*[6pt]
\cline{2-5}\noalign{\nobreak}
 & Design Moment Capacity, $M_d$ & IRC 22 Cl. 603.3.1 & 11371.85 kN-m & --- \\*[6pt]
\cline{2-5}\noalign{\nobreak}
 & Utilization Ratio, $M_u / M_d$ & $M_u / M_d$ & 0.20 & PASS \\*[6pt]%
\end{longtable}
\noindent\textit{Note: IRC 22 Cl. 603.3.1, IS 800 Cl. 8.2.1}

\vspace{1em}
\begin{longtable}{|C{2.20cm}|C{3.08cm}|C{3.08cm}|>{\centering\arraybackslash}p{3.70cm}|C{1.58cm}|}%
\caption{\textbf{Shear Capacity Check}}\\%
\hline%
\rowcolor{tableHeaderBg}\hspace{0pt}\textbf{} & \hspace{0pt}\textbf{Parameter} & \hspace{0pt}\textbf{Formula} & \hspace{0pt}\textbf{Value} & \hspace{0pt}\textbf{Status} \\%
\hline%
\endfirsthead%
\caption[]{\textbf{Shear Capacity Check} (continued)}\\%
\hline%
\rowcolor{tableHeaderBg}\hspace{0pt}\textbf{} & \hspace{0pt}\textbf{Parameter} & \hspace{0pt}\textbf{Formula} & \hspace{0pt}\textbf{Value} & \hspace{0pt}\textbf{Status} \\%
\hline%
\endhead%
\multicolumn{5}{r}{{\footnotesize\itshape Continued on next page}} \\%
\endfoot%
\hline%
\endlastfoot%
\multirow{8}{*}{\makecell{G1}} & Applied Shear, $V_u$ & Governing LC (ULS) & 288.29 kN & --- \\*[6pt]
\cline{2-5}\noalign{\nobreak}
 & Shear Area, $A_v$ & $d_w \times t_w$ & 16260.0 mm$^2$ & --- \\*[6pt]
\cline{2-5}\noalign{\nobreak}
 & Panel Aspect Ratio, c/d & --- & 1.501 & --- \\*[6pt]
\cline{2-5}\noalign{\nobreak}
 & Shear Buckling Coefficient, $k_v$ & IS 800 Cl. 8.4.2.2 & 7.126 & --- \\*[6pt]
\cline{2-5}\noalign{\nobreak}
 & Web Slenderness, $\lambda_w$ & $\sqrt{f_{yw}/(\sqrt{3}\,\tau_{cr})}$ & 2.037 & --- \\*[6pt]
\cline{2-5}\noalign{\nobreak}
 & Design Shear Stress, $\tau_b$ & IRC 22 Cl. 603.3.3.2 & 48.72 MPa & --- \\*[6pt]
\cline{2-5}\noalign{\nobreak}
 & Shear Buckling Resistance, $V_{cr}$ & $A_v \times \tau_b$ & 720.21 kN & --- \\*[6pt]
\cline{2-5}\noalign{\nobreak}
 & Utilization Ratio, $V_u / V_d$ & $V_u / V_d$ & 0.40 & PASS \\*[6pt]
\hline
\multirow{8}{*}{\makecell{G2}} & Applied Shear, $V_u$ & Governing LC (ULS) & 288.29 kN & --- \\*[6pt]
\cline{2-5}\noalign{\nobreak}
 & Shear Area, $A_v$ & $d_w \times t_w$ & 16260.0 mm$^2$ & --- \\*[6pt]
\cline{2-5}\noalign{\nobreak}
 & Panel Aspect Ratio, c/d & --- & 1.501 & --- \\*[6pt]
\cline{2-5}\noalign{\nobreak}
 & Shear Buckling Coefficient, $k_v$ & IS 800 Cl. 8.4.2.2 & 7.126 & --- \\*[6pt]
\cline{2-5}\noalign{\nobreak}
 & Web Slenderness, $\lambda_w$ & $\sqrt{f_{yw}/(\sqrt{3}\,\tau_{cr})}$ & 2.037 & --- \\*[6pt]
\cline{2-5}\noalign{\nobreak}
 & Design Shear Stress, $\tau_b$ & IRC 22 Cl. 603.3.3.2 & 48.72 MPa & --- \\*[6pt]
\cline{2-5}\noalign{\nobreak}
 & Shear Buckling Resistance, $V_{cr}$ & $A_v \times \tau_b$ & 720.21 kN & --- \\*[6pt]
\cline{2-5}\noalign{\nobreak}
 & Utilization Ratio, $V_u / V_d$ & $V_u / V_d$ & 0.40 & PASS \\*[6pt]
\hline
\multirow{8}{*}{\makecell{G3}} & Applied Shear, $V_u$ & Governing LC (ULS) & 288.29 kN & --- \\*[6pt]
\cline{2-5}\noalign{\nobreak}
 & Shear Area, $A_v$ & $d_w \times t_w$ & 16260.0 mm$^2$ & --- \\*[6pt]
\cline{2-5}\noalign{\nobreak}
 & Panel Aspect Ratio, c/d & --- & 1.501 & --- \\*[6pt]
\cline{2-5}\noalign{\nobreak}
 & Shear Buckling Coefficient, $k_v$ & IS 800 Cl. 8.4.2.2 & 7.126 & --- \\*[6pt]
\cline{2-5}\noalign{\nobreak}
 & Web Slenderness, $\lambda_w$ & $\sqrt{f_{yw}/(\sqrt{3}\,\tau_{cr})}$ & 2.037 & --- \\*[6pt]
\cline{2-5}\noalign{\nobreak}
 & Design Shear Stress, $\tau_b$ & IRC 22 Cl. 603.3.3.2 & 48.72 MPa & --- \\*[6pt]
\cline{2-5}\noalign{\nobreak}
 & Shear Buckling Resistance, $V_{cr}$ & $A_v \times \tau_b$ & 720.21 kN & --- \\*[6pt]
\cline{2-5}\noalign{\nobreak}
 & Utilization Ratio, $V_u / V_d$ & $V_u / V_d$ & 0.40 & PASS \\*[6pt]
\hline
\multirow{8}{*}{\makecell{G4}} & Applied Shear, $V_u$ & Governing LC (ULS) & 288.29 kN & --- \\*[6pt]
\cline{2-5}\noalign{\nobreak}
 & Shear Area, $A_v$ & $d_w \times t_w$ & 16260.0 mm$^2$ & --- \\*[6pt]
\cline{2-5}\noalign{\nobreak}
 & Panel Aspect Ratio, c/d & --- & 1.501 & --- \\*[6pt]
\cline{2-5}\noalign{\nobreak}
 & Shear Buckling Coefficient, $k_v$ & IS 800 Cl. 8.4.2.2 & 7.126 & --- \\*[6pt]
\cline{2-5}\noalign{\nobreak}
 & Web Slenderness, $\lambda_w$ & $\sqrt{f_{yw}/(\sqrt{3}\,\tau_{cr})}$ & 2.037 & --- \\*[6pt]
\cline{2-5}\noalign{\nobreak}
 & Design Shear Stress, $\tau_b$ & IRC 22 Cl. 603.3.3.2 & 48.72 MPa & --- \\*[6pt]
\cline{2-5}\noalign{\nobreak}
 & Shear Buckling Resistance, $V_{cr}$ & $A_v \times \tau_b$ & 720.21 kN & --- \\*[6pt]
\cline{2-5}\noalign{\nobreak}
 & Utilization Ratio, $V_u / V_d$ & $V_u / V_d$ & 0.40 & PASS \\*[6pt]%
\end{longtable}
\noindent\textit{Note: IS 800 Cl. 8.4, IRC 22 Cl. 603.3.3.2}

\vspace{1em}
\begin{longtable}{|C{2.20cm}|C{3.08cm}|C{3.08cm}|>{\centering\arraybackslash}p{3.70cm}|C{1.58cm}|}%
\caption{\textbf{Interaction Checks (M-V and M-N)}}\\%
\hline%
\rowcolor{tableHeaderBg}\hspace{0pt}\textbf{} & \hspace{0pt}\textbf{Check} & \hspace{0pt}\textbf{Condition} & \hspace{0pt}\textbf{Value} & \hspace{0pt}\textbf{Status} \\%
\hline%
\endfirsthead%
\caption[]{\textbf{Interaction Checks (M-V and M-N)} (continued)}\\%
\hline%
\rowcolor{tableHeaderBg}\hspace{0pt}\textbf{} & \hspace{0pt}\textbf{Check} & \hspace{0pt}\textbf{Condition} & \hspace{0pt}\textbf{Value} & \hspace{0pt}\textbf{Status} \\%
\hline%
\endhead%
\multicolumn{5}{r}{{\footnotesize\itshape Continued on next page}} \\%
\endfoot%
\hline%
\endlastfoot%
\multirow{4}{*}{\makecell{G1}} & High Shear Condition? & $V_u > 0.6\,V_d$ & No & --- \\*[6pt]
\cline{2-5}\noalign{\nobreak}
 & Reduced Moment Capacity, $M_{dv}$ & IRC 22 Cl. 603.3.3.3 & 11371.85 kN-m & --- \\*[6pt]
\cline{2-5}\noalign{\nobreak}
 & Interaction Check: $M_u \leq M_{dv}$ & --- & 0.20 & PASS \\*[6pt]
\cline{2-5}\noalign{\nobreak}
 & Interaction Check: $N_u/N_{Rd} + M_u/M_{dv} \leq 1.0$ & 0.00 + 0.20 = 0.198 & 0.198 & PASS \\*[6pt]
\hline
\multirow{4}{*}{\makecell{G2}} & High Shear Condition? & $V_u > 0.6\,V_d$ & No & --- \\*[6pt]
\cline{2-5}\noalign{\nobreak}
 & Reduced Moment Capacity, $M_{dv}$ & IRC 22 Cl. 603.3.3.3 & 11371.85 kN-m & --- \\*[6pt]
\cline{2-5}\noalign{\nobreak}
 & Interaction Check: $M_u \leq M_{dv}$ & --- & 0.20 & PASS \\*[6pt]
\cline{2-5}\noalign{\nobreak}
 & Interaction Check: $N_u/N_{Rd} + M_u/M_{dv} \leq 1.0$ & 0.00 + 0.20 = 0.198 & 0.198 & PASS \\*[6pt]
\hline
\multirow{4}{*}{\makecell{G3}} & High Shear Condition? & $V_u > 0.6\,V_d$ & No & --- \\*[6pt]
\cline{2-5}\noalign{\nobreak}
 & Reduced Moment Capacity, $M_{dv}$ & IRC 22 Cl. 603.3.3.3 & 11371.85 kN-m & --- \\*[6pt]
\cline{2-5}\noalign{\nobreak}
 & Interaction Check: $M_u \leq M_{dv}$ & --- & 0.20 & PASS \\*[6pt]
\cline{2-5}\noalign{\nobreak}
 & Interaction Check: $N_u/N_{Rd} + M_u/M_{dv} \leq 1.0$ & 0.00 + 0.20 = 0.198 & 0.198 & PASS \\*[6pt]
\hline
\multirow{4}{*}{\makecell{G4}} & High Shear Condition? & $V_u > 0.6\,V_d$ & No & --- \\*[6pt]
\cline{2-5}\noalign{\nobreak}
 & Reduced Moment Capacity, $M_{dv}$ & IRC 22 Cl. 603.3.3.3 & 11371.85 kN-m & --- \\*[6pt]
\cline{2-5}\noalign{\nobreak}
 & Interaction Check: $M_u \leq M_{dv}$ & --- & 0.20 & PASS \\*[6pt]
\cline{2-5}\noalign{\nobreak}
 & Interaction Check: $N_u/N_{Rd} + M_u/M_{dv} \leq 1.0$ & 0.00 + 0.20 = 0.198 & 0.198 & PASS \\*[6pt]%
\end{longtable}
\noindent\textit{Note: IRC 22 Cl. 603.3.3.3}

\vspace{1em}
\begin{longtable}{|C{2.20cm}|C{3.08cm}|C{3.08cm}|>{\centering\arraybackslash}p{3.70cm}|C{1.58cm}|}%
\caption{\textbf{Lateral Torsional Buckling Check -- Construction Stage}}\\%
\hline%
\rowcolor{tableHeaderBg}\hspace{0pt}\textbf{} & \hspace{0pt}\textbf{Parameter} & \hspace{0pt}\textbf{Formula} & \hspace{0pt}\textbf{Value} & \hspace{0pt}\textbf{Status} \\%
\hline%
\endfirsthead%
\caption[]{\textbf{Lateral Torsional Buckling Check -- Construction Stage} (continued)}\\%
\hline%
\rowcolor{tableHeaderBg}\hspace{0pt}\textbf{} & \hspace{0pt}\textbf{Parameter} & \hspace{0pt}\textbf{Formula} & \hspace{0pt}\textbf{Value} & \hspace{0pt}\textbf{Status} \\%
\hline%
\endhead%
\multicolumn{5}{r}{{\footnotesize\itshape Continued on next page}} \\%
\endfoot%
\hline%
\endlastfoot%
\multirow{5}{*}{\makecell{G1}} & Elastic Critical Moment, $M_{cr}$ & IRC 22 Cl. 603.3.3.1 & 88107.8 kN-m & --- \\*[6pt]
\cline{2-5}\noalign{\nobreak}
 & Non-dim. Slenderness, $\bar{\lambda}_{LT}$ & $\sqrt{M_p / M_{cr}}$ & 0.299 & --- \\*[6pt]
\cline{2-5}\noalign{\nobreak}
 & LTB Reduction Factor, $\chi_{LT}$ & IS 800 Cl. 8.2.2 & 1.0 & --- \\*[6pt]
\cline{2-5}\noalign{\nobreak}
 & LTB Resistance, $M_b$ & $\chi_{LT}\,M_p / \gamma_{m0}$ & 7986.44 kN-m & --- \\*[6pt]
\cline{2-5}\noalign{\nobreak}
 & $M_u \leq M_b$ & $M_u / M_b$ & 0.23 & PASS \\*[6pt]
\hline
\multirow{5}{*}{\makecell{G2}} & Elastic Critical Moment, $M_{cr}$ & IRC 22 Cl. 603.3.3.1 & 88107.8 kN-m & --- \\*[6pt]
\cline{2-5}\noalign{\nobreak}
 & Non-dim. Slenderness, $\bar{\lambda}_{LT}$ & $\sqrt{M_p / M_{cr}}$ & 0.299 & --- \\*[6pt]
\cline{2-5}\noalign{\nobreak}
 & LTB Reduction Factor, $\chi_{LT}$ & IS 800 Cl. 8.2.2 & 1.0 & --- \\*[6pt]
\cline{2-5}\noalign{\nobreak}
 & LTB Resistance, $M_b$ & $\chi_{LT}\,M_p / \gamma_{m0}$ & 7986.44 kN-m & --- \\*[6pt]
\cline{2-5}\noalign{\nobreak}
 & $M_u \leq M_b$ & $M_u / M_b$ & 0.23 & PASS \\*[6pt]
\hline
\multirow{5}{*}{\makecell{G3}} & Elastic Critical Moment, $M_{cr}$ & IRC 22 Cl. 603.3.3.1 & 88107.8 kN-m & --- \\*[6pt]
\cline{2-5}\noalign{\nobreak}
 & Non-dim. Slenderness, $\bar{\lambda}_{LT}$ & $\sqrt{M_p / M_{cr}}$ & 0.299 & --- \\*[6pt]
\cline{2-5}\noalign{\nobreak}
 & LTB Reduction Factor, $\chi_{LT}$ & IS 800 Cl. 8.2.2 & 1.0 & --- \\*[6pt]
\cline{2-5}\noalign{\nobreak}
 & LTB Resistance, $M_b$ & $\chi_{LT}\,M_p / \gamma_{m0}$ & 7986.44 kN-m & --- \\*[6pt]
\cline{2-5}\noalign{\nobreak}
 & $M_u \leq M_b$ & $M_u / M_b$ & 0.23 & PASS \\*[6pt]
\hline
\multirow{5}{*}{\makecell{G4}} & Elastic Critical Moment, $M_{cr}$ & IRC 22 Cl. 603.3.3.1 & 88107.8 kN-m & --- \\*[6pt]
\cline{2-5}\noalign{\nobreak}
 & Non-dim. Slenderness, $\bar{\lambda}_{LT}$ & $\sqrt{M_p / M_{cr}}$ & 0.299 & --- \\*[6pt]
\cline{2-5}\noalign{\nobreak}
 & LTB Reduction Factor, $\chi_{LT}$ & IS 800 Cl. 8.2.2 & 1.0 & --- \\*[6pt]
\cline{2-5}\noalign{\nobreak}
 & LTB Resistance, $M_b$ & $\chi_{LT}\,M_p / \gamma_{m0}$ & 7986.44 kN-m & --- \\*[6pt]
\cline{2-5}\noalign{\nobreak}
 & $M_u \leq M_b$ & $M_u / M_b$ & 0.23 & PASS \\*[6pt]%
\end{longtable}
\noindent\textit{Note: IRC 22 Cl. 603.3.3.1, IS 800 Cl. 8.2.2}

\vspace{1em}
\begin{longtable}{|C{2.34cm}|L{6.09cm}|>{\arraybackslash}p{6.09cm}|}%
\caption{\textbf{Stiffener Design Summary}}\\%
\hline%
\endfirsthead%
\caption[]{\textbf{Stiffener Design Summary} (continued)}\\%
\hline%
\endhead%
\multicolumn{3}{r}{{\footnotesize\itshape Continued on next page}} \\%
\endfoot%
\hline%
\endlastfoot%
\multirow{6}{*}{\makecell{G1}} & \textnormal{Shear Buckling Design Method} & post\_\allowbreak{}critical \\*[6pt]
\cline{2-3}\noalign{\nobreak}
 & \textnormal{Required Intermediate Stiffener Thickness (mm)} & 21.1 \\*[6pt]
\cline{2-3}\noalign{\nobreak}
 & \textnormal{Required Intermediate Stiffener Spacing (mm)} & \textit{Not available} \\*[6pt]
\cline{2-3}\noalign{\nobreak}
 & \textnormal{Required End Panel Stiffener Thickness (mm)} & 1.2 \\*[6pt]
\cline{2-3}\noalign{\nobreak}
 & \textnormal{No. of End Panel Stiffeners} & 4 \\*[6pt]
\cline{2-3}\noalign{\nobreak}
 & \textnormal{Longitudinal Stiffeners} & None \\*[6pt]
\hline
\multirow{6}{*}{\makecell{G2}} & \textnormal{Shear Buckling Design Method} & post\_\allowbreak{}critical \\*[6pt]
\cline{2-3}\noalign{\nobreak}
 & \textnormal{Required Intermediate Stiffener Thickness (mm)} & 21.1 \\*[6pt]
\cline{2-3}\noalign{\nobreak}
 & \textnormal{Required Intermediate Stiffener Spacing (mm)} & \textit{Not available} \\*[6pt]
\cline{2-3}\noalign{\nobreak}
 & \textnormal{Required End Panel Stiffener Thickness (mm)} & 1.2 \\*[6pt]
\cline{2-3}\noalign{\nobreak}
 & \textnormal{No. of End Panel Stiffeners} & 4 \\*[6pt]
\cline{2-3}\noalign{\nobreak}
 & \textnormal{Longitudinal Stiffeners} & None \\*[6pt]
\hline
\multirow{6}{*}{\makecell{G3}} & \textnormal{Shear Buckling Design Method} & post\_\allowbreak{}critical \\*[6pt]
\cline{2-3}\noalign{\nobreak}
 & \textnormal{Required Intermediate Stiffener Thickness (mm)} & 21.1 \\*[6pt]
\cline{2-3}\noalign{\nobreak}
 & \textnormal{Required Intermediate Stiffener Spacing (mm)} & \textit{Not available} \\*[6pt]
\cline{2-3}\noalign{\nobreak}
 & \textnormal{Required End Panel Stiffener Thickness (mm)} & 1.2 \\*[6pt]
\cline{2-3}\noalign{\nobreak}
 & \textnormal{No. of End Panel Stiffeners} & 4 \\*[6pt]
\cline{2-3}\noalign{\nobreak}
 & \textnormal{Longitudinal Stiffeners} & None \\*[6pt]
\hline
\multirow{6}{*}{\makecell{G4}} & \textnormal{Shear Buckling Design Method} & post\_\allowbreak{}critical \\*[6pt]
\cline{2-3}\noalign{\nobreak}
 & \textnormal{Required Intermediate Stiffener Thickness (mm)} & 21.1 \\*[6pt]
\cline{2-3}\noalign{\nobreak}
 & \textnormal{Required Intermediate Stiffener Spacing (mm)} & \textit{Not available} \\*[6pt]
\cline{2-3}\noalign{\nobreak}
 & \textnormal{Required End Panel Stiffener Thickness (mm)} & 1.2 \\*[6pt]
\cline{2-3}\noalign{\nobreak}
 & \textnormal{No. of End Panel Stiffeners} & 4 \\*[6pt]
\cline{2-3}\noalign{\nobreak}
 & \textnormal{Longitudinal Stiffeners} & None \\*[6pt]%
\end{longtable}

\vspace{1em}

\begin{longtable}{|C{2.20cm}|L{3.08cm}|C{3.08cm}|>{\centering\arraybackslash}p{3.70cm}|C{1.58cm}|}%
\caption{\textbf{End Panel Stiffener Checks}}\\%
\hline%
\rowcolor{tableHeaderBg}\hspace{0pt}\textbf{} & \hspace{0pt}\textbf{Check} & \hspace{0pt}\textbf{Required} & \hspace{0pt}\textbf{Provided} & \hspace{0pt}\textbf{Status} \\%
\hline%
\endfirsthead%
\caption[]{\textbf{End Panel Stiffener Checks} (continued)}\\%
\hline%
\rowcolor{tableHeaderBg}\hspace{0pt}\textbf{} & \hspace{0pt}\textbf{Check} & \hspace{0pt}\textbf{Required} & \hspace{0pt}\textbf{Provided} & \hspace{0pt}\textbf{Status} \\%
\hline%
\endhead%
\multicolumn{5}{r}{{\footnotesize\itshape Continued on next page}} \\%
\endfoot%
\hline%
\endlastfoot%
\multirow{4}{*}{\makecell{G1}} & Web Buckling Resistance & 379.98 kN & 168429.67 kN & PASS \\*[6pt]
\cline{2-5}\noalign{\nobreak}
 & Local Crushing Resistance & 379.98 kN & 484.98 kN & PASS \\*[6pt]
\cline{2-5}\noalign{\nobreak}
 & Bearing Capacity & 379.98 kN & 3181.82 kN & PASS \\*[6pt]
\cline{2-5}\noalign{\nobreak}
 & Column Buckling Resistance & 379.98 kN & 2863.64 kN & PASS \\*[6pt]
\hline
\multirow{4}{*}{\makecell{G2}} & Web Buckling Resistance & 379.98 kN & 168429.67 kN & PASS \\*[6pt]
\cline{2-5}\noalign{\nobreak}
 & Local Crushing Resistance & 379.98 kN & 484.98 kN & PASS \\*[6pt]
\cline{2-5}\noalign{\nobreak}
 & Bearing Capacity & 379.98 kN & 3181.82 kN & PASS \\*[6pt]
\cline{2-5}\noalign{\nobreak}
 & Column Buckling Resistance & 379.98 kN & 2863.64 kN & PASS \\*[6pt]
\hline
\multirow{4}{*}{\makecell{G3}} & Web Buckling Resistance & 379.98 kN & 168429.67 kN & PASS \\*[6pt]
\cline{2-5}\noalign{\nobreak}
 & Local Crushing Resistance & 379.98 kN & 484.98 kN & PASS \\*[6pt]
\cline{2-5}\noalign{\nobreak}
 & Bearing Capacity & 379.98 kN & 3181.82 kN & PASS \\*[6pt]
\cline{2-5}\noalign{\nobreak}
 & Column Buckling Resistance & 379.98 kN & 2863.64 kN & PASS \\*[6pt]
\hline
\multirow{4}{*}{\makecell{G4}} & Web Buckling Resistance & 379.98 kN & 168429.67 kN & PASS \\*[6pt]
\cline{2-5}\noalign{\nobreak}
 & Local Crushing Resistance & 379.98 kN & 484.98 kN & PASS \\*[6pt]
\cline{2-5}\noalign{\nobreak}
 & Bearing Capacity & 379.98 kN & 3181.82 kN & PASS \\*[6pt]
\cline{2-5}\noalign{\nobreak}
 & Column Buckling Resistance & 379.98 kN & 2863.64 kN & PASS \\*[6pt]%
\end{longtable}
\noindent\textit{Note: IS 800 Cl. 8.4.2.2}

\vspace{1em}
\begin{longtable}{|C{2.20cm}|L{3.08cm}|C{3.08cm}|>{\centering\arraybackslash}p{3.08cm}|C{2.20cm}|}%
\caption{\textbf{Serviceability -- Deflection Checks}}\\%
\hline%
\rowcolor{tableHeaderBg}\hspace{0pt}\textbf{} & \hspace{0pt}\textbf{Check} & \hspace{0pt}\textbf{Allowable} & \hspace{0pt}\textbf{Actual} & \hspace{0pt}\textbf{Status} \\%
\hline%
\endfirsthead%
\caption[]{\textbf{Serviceability -- Deflection Checks} (continued)}\\%
\hline%
\rowcolor{tableHeaderBg}\hspace{0pt}\textbf{} & \hspace{0pt}\textbf{Check} & \hspace{0pt}\textbf{Allowable} & \hspace{0pt}\textbf{Actual} & \hspace{0pt}\textbf{Status} \\%
\hline%
\endhead%
\multicolumn{5}{r}{{\footnotesize\itshape Continued on next page}} \\%
\endfoot%
\hline%
\endlastfoot%
\multirow{2}{*}{\makecell{G1}} & Live Load Deflection, $\delta_{LL}$ (mm) & L/800 = 37.5 mm & 19.252 & PASS \\*[6pt]
\cline{2-5}\noalign{\nobreak}
 & Total Load Deflection, $\delta_{total}$ (mm) & L/600 = 50.0 mm & 51.090 & \textcolor{red}{FAIL} \\*[6pt]
\hline
\multirow{2}{*}{\makecell{G2}} & Live Load Deflection, $\delta_{LL}$ (mm) & L/800 = 37.5 mm & 19.385 & PASS \\*[6pt]
\cline{2-5}\noalign{\nobreak}
 & Total Load Deflection, $\delta_{total}$ (mm) & L/600 = 50.0 mm & 51.225 & \textcolor{red}{FAIL} \\*[6pt]
\hline
\multirow{2}{*}{\makecell{G3}} & Live Load Deflection, $\delta_{LL}$ (mm) & L/800 = 37.5 mm & 19.385 & PASS \\*[6pt]
\cline{2-5}\noalign{\nobreak}
 & Total Load Deflection, $\delta_{total}$ (mm) & L/600 = 50.0 mm & 51.225 & \textcolor{red}{FAIL} \\*[6pt]
\hline
\multirow{2}{*}{\makecell{G4}} & Live Load Deflection, $\delta_{LL}$ (mm) & L/800 = 37.5 mm & 19.252 & PASS \\*[6pt]
\cline{2-5}\noalign{\nobreak}
 & Total Load Deflection, $\delta_{total}$ (mm) & L/600 = 50.0 mm & 51.090 & \textcolor{red}{FAIL} \\*[6pt]%
\end{longtable}
\noindent\textit{Note: IRC 22 Cl. 604.3.2}

\vspace{1em}
\begin{longtable}{|C{2.20cm}|L{3.08cm}|C{3.08cm}|>{\centering\arraybackslash}p{3.08cm}|C{2.20cm}|}%
\caption{\textbf{Serviceability -- Maximum Stress Limitation}}\\%
\hline%
\rowcolor{tableHeaderBg}\hspace{0pt}\textbf{} & \hspace{0pt}\textbf{Element} & \hspace{0pt}\textbf{Allowable Stress} & \hspace{0pt}\textbf{Actual Stress} & \hspace{0pt}\textbf{Status} \\%
\hline%
\endfirsthead%
\caption[]{\textbf{Serviceability -- Maximum Stress Limitation} (continued)}\\%
\hline%
\rowcolor{tableHeaderBg}\hspace{0pt}\textbf{} & \hspace{0pt}\textbf{Element} & \hspace{0pt}\textbf{Allowable Stress} & \hspace{0pt}\textbf{Actual Stress} & \hspace{0pt}\textbf{Status} \\%
\hline%
\endhead%
\multicolumn{5}{r}{{\footnotesize\itshape Continued on next page}} \\%
\endfoot%
\hline%
\endlastfoot%
\makecell{G1} & Structural Steel ($0.9\,f_y$) & 315.00 MPa & 78.56 MPa & PASS \\*[6pt]
\hline
\makecell{G2} & Structural Steel ($0.9\,f_y$) & 315.00 MPa & 78.56 MPa & PASS \\*[6pt]
\hline
\makecell{G3} & Structural Steel ($0.9\,f_y$) & 315.00 MPa & 78.56 MPa & PASS \\*[6pt]
\hline
\makecell{G4} & Structural Steel ($0.9\,f_y$) & 315.00 MPa & 78.56 MPa & PASS \\*[6pt]%
\end{longtable}

\vspace{1em}
\begin{longtable}{|C{2.20cm}|C{3.08cm}|C{3.08cm}|>{\centering\arraybackslash}p{3.08cm}|C{2.20cm}|}%
\caption{\textbf{Serviceability -- Fatigue Assessment}}\\%
\hline%
\rowcolor{tableHeaderBg}\hspace{0pt}\textbf{} & \hspace{0pt}\textbf{Stress Range, $\Delta\sigma$ (MPa)} & \hspace{0pt}\textbf{Fatigue Limit, $f_{fd}$ (MPa)} & \hspace{0pt}\textbf{Utilization Ratio} & \hspace{0pt}\textbf{Status} \\%
\hline%
\endfirsthead%
\caption[]{\textbf{Serviceability -- Fatigue Assessment} (continued)}\\%
\hline%
\rowcolor{tableHeaderBg}\hspace{0pt}\textbf{} & \hspace{0pt}\textbf{Stress Range, $\Delta\sigma$ (MPa)} & \hspace{0pt}\textbf{Fatigue Limit, $f_{fd}$ (MPa)} & \hspace{0pt}\textbf{Utilization Ratio} & \hspace{0pt}\textbf{Status} \\%
\hline%
\endhead%
\multicolumn{5}{r}{{\footnotesize\itshape Continued on next page}} \\%
\endfoot%
\hline%
\endlastfoot%
\makecell{G1} & 142.19 MPa & 92.49 MPa & 1.54 & \textcolor{red}{FAIL} \\*[6pt]
\hline
\makecell{G2} & 142.86 MPa & 92.49 MPa & 1.54 & \textcolor{red}{FAIL} \\*[6pt]
\hline
\makecell{G3} & 145.86 MPa & 92.49 MPa & 1.58 & \textcolor{red}{FAIL} \\*[6pt]
\hline
\makecell{G4} & 145.86 MPa & 92.49 MPa & 1.58 & \textcolor{red}{FAIL} \\*[6pt]%
\end{longtable}
\noindent\textit{Note: IRC 22 Cl. 605 --- governing of normal and shear fatigue (worst by DCR). Capacity reduction factor $\mu_r$ applied where plate thickness > 25 mm.}

\vspace{1em}
\vspace{0.4em}
\begin{longtable}{|C{1.37cm}|>{\centering\arraybackslash}p{3.08cm}|C{2.05cm}|C{1.71cm}|C{1.80cm}|C{1.45cm}|C{1.28cm}|}%
\caption{\textbf{Girder Design Summary (DCR / Utilization Ratio)}}\\%
\hline%
\rowcolor{tableHeaderBg}\hspace{0pt}\textbf{Girder} & \hspace{0pt}\textbf{Controlling LC / Combination} & \hspace{0pt}\textbf{Controlling Check} & \hspace{0pt}\textbf{Demand} & \hspace{0pt}\textbf{Capacity} & \hspace{0pt}\textbf{UR} & \hspace{0pt}\textbf{Status} \\%
\hline%
\endfirsthead%
\caption[]{\textbf{Girder Design Summary (DCR / Utilization Ratio)} (continued)}\\%
\hline%
\rowcolor{tableHeaderBg}\hspace{0pt}\textbf{Girder} & \hspace{0pt}\textbf{Controlling LC / Combination} & \hspace{0pt}\textbf{Controlling Check} & \hspace{0pt}\textbf{Demand} & \hspace{0pt}\textbf{Capacity} & \hspace{0pt}\textbf{UR} & \hspace{0pt}\textbf{Status} \\%
\hline%
\endhead%
\multicolumn{7}{r}{{\footnotesize\itshape Continued on next page}} \\%
\endfoot%
\hline%
\endlastfoot%
G1 & SLS\_\allowbreak{}FREQUENT\_\allowbreak{}1:\allowbreak{} 1.0DL + 1.2DW + 0.75LL + 0.5WL & Fatigue Normal Stress & 142.19 MPa & 92.49 MPa & 1.537 & \textcolor{red}{FAIL} \\*[6pt]
\hline
G2 & SLS\_\allowbreak{}FREQUENT\_\allowbreak{}1:\allowbreak{} 1.0DL + 1.2DW + 0.75LL + 0.5WL & Fatigue Normal Stress & 142.86 MPa & 92.49 MPa & 1.545 & \textcolor{red}{FAIL} \\*[6pt]
\hline
G3 & SLS\_\allowbreak{}FREQUENT\_\allowbreak{}1:\allowbreak{} 1.0DL + 1.2DW + 0.75LL + 0.5WL & Fatigue Normal Stress & 145.86 MPa & 92.49 MPa & 1.577 & \textcolor{red}{FAIL} \\*[6pt]
\hline
G4 & SLS\_\allowbreak{}FREQUENT\_\allowbreak{}1:\allowbreak{} 1.0DL + 1.2DW + 0.75LL + 0.5WL & Fatigue Normal Stress & 145.86 MPa & 92.49 MPa & 1.577 & \textcolor{red}{FAIL} \\*[6pt]%
\end{longtable}
\noindent\textit{Note: UR = Demand / Capacity. A value $\leq 1.0$ indicates a passing check. The controlling check is the criterion with the highest UR for each girder, with the real load case/combination that drives it.}

\vspace{1em}

\begin{longtable}{|L{3.42cm}|C{5.33cm}|>{\centering\arraybackslash}p{2.47cm}|L{2.85cm}|}%
\caption{\textbf{Shear Connector Capacity}}\\%
\hline%
\rowcolor{tableHeaderBg}\hspace{0pt}\textbf{Parameter} & \hspace{0pt}\textbf{Formula} & \hspace{0pt}\textbf{Value} & \hspace{0pt}\textbf{Reference} \\%
\hline%
\endfirsthead%
\caption[]{\textbf{Shear Connector Capacity} (continued)}\\%
\hline%
\rowcolor{tableHeaderBg}\hspace{0pt}\textbf{Parameter} & \hspace{0pt}\textbf{Formula} & \hspace{0pt}\textbf{Value} & \hspace{0pt}\textbf{Reference} \\%
\hline%
\endhead%
\multicolumn{4}{r}{{\footnotesize\itshape Continued on next page}} \\%
\endfoot%
\hline%
\endlastfoot%
Design Resistance, $Q_u$ & \footnotesize\makecell{$Q_u=\min(Q_{u,s},\,Q_{u,c})$\\[3pt]$Q_{u,s}=\dfrac{0.8\,f_u\,(\pi d^2/4)}{\gamma_v}$\\[3pt]$Q_{u,c}=\dfrac{0.29\,\alpha\,d^2\sqrt{f_{ck}\,E_{cm}}}{\gamma_v}$} & 95.36 kN & IRC 22 Cl. 606.3.1 (Eq. 6.1) \\*[6pt]
\hline
Fatigue Shear Resistance, $Q_r$ & IRC 22 Table 8 ($\phi d$, $N_{sc}$) & 25.00 kN & IRC 22 Cl. 606.3.2 (Table 8) \\*[6pt]%
\end{longtable}

\vspace{1em}


\begin{longtable}{|L{3.2cm}|>{\centering\arraybackslash}p{4.3cm}|>{\centering\arraybackslash}p{4.3cm}|C{2.0cm}|}%
\caption{\textbf{Shear Connector Spacing}}\\%
\hline%
\rowcolor{tableHeaderBg}\hspace{0pt}\textbf{Criterion} & \hspace{0pt}\textbf{Governing Spacing} & \hspace{0pt}\textbf{Actual Spacing Provided} & \hspace{0pt}\textbf{Status} \\%
\hline%
\endfirsthead%
\caption[]{\textbf{Shear Connector Spacing} (continued)}\\%
\hline%
\rowcolor{tableHeaderBg}\hspace{0pt}\textbf{Criterion} & \hspace{0pt}\textbf{Governing Spacing} & \hspace{0pt}\textbf{Actual Spacing Provided} & \hspace{0pt}\textbf{Status} \\%
\hline%
\endhead%
\multicolumn{4}{r}{{\footnotesize\itshape Continued on next page}} \\%
\endfoot%
\hline%
\endlastfoot%
ULS Shear (SL1) & 2302.4 mm & 378.4 mm & PASS \\*[6pt]
\hline
Full Composite (SL2) & 378.4 mm & 378.4 mm & PASS \\*[6pt]
\hline
SLS Fatigue (SR) & 665.4 mm & 378.4 mm & PASS \\*[6pt]
\hline
Max Spacing Limit (IRC 22) & 400.0 mm & 378.4 mm & PASS \\*[6pt]%
\end{longtable}
\noindent\textit{Note: IRC 22 Cl. 606.4, 606.9. Governing spacing $= \min(S_{L1}, S_{L2}, S_R)$.}

\vspace{1em}
\begin{longtable}{|L{5.3cm}|>{\arraybackslash}p{7.2cm}|C{2.0cm}|}%
\caption{\textbf{Transverse Shear and Detailing Checks}}\\%
\hline%
\rowcolor{tableHeaderBg}\hspace{0pt}\textbf{Check} & \hspace{0pt}\textbf{Value} & \hspace{0pt}\textbf{Status} \\%
\hline%
\endfirsthead%
\caption[]{\textbf{Transverse Shear and Detailing Checks} (continued)}\\%
\hline%
\rowcolor{tableHeaderBg}\hspace{0pt}\textbf{Check} & \hspace{0pt}\textbf{Value} & \hspace{0pt}\textbf{Status} \\%
\hline%
\endhead%
\multicolumn{3}{r}{{\footnotesize\itshape Continued on next page}} \\%
\endfoot%
\hline%
\endlastfoot%
\textnormal{Longitudinal Shear per unit length, $V_L$} & 82.84 kN/m & --- \\*[6pt]
\hline
\textnormal{Transverse Shear Capacity of Slab, $V_{Rd}$} & 999.28 kN/m & --- \\*[6pt]
\hline
\textnormal{Transverse Shear Check} & $V_L/V_{Rd}$ = 0.08 & PASS \\*[6pt]
\hline
\textnormal{Min. Transverse Reinforcement, $A_{st,min}$} & Required 0.41 cm$^2$/m, Provided 11.38 cm$^2$/m & PASS \\*[6pt]
\hline
\textnormal{Stud Diameter $\leq 2\,t_f$} & $d$ = 20.0 mm $\leq 2t_f$ = 44.0 mm & PASS \\*[6pt]
\hline
\textnormal{Stud Edge Distance} & Provided 195.0 mm (req. $\geq$ 25.0 mm) & PASS \\*[6pt]%
\end{longtable}
\noindent\textit{Note: IRC 22 Cl. 606.6, 606.10.}

% ===========================
\section{Deck Slab Design}
\label{sec:deck-design}
% ===========================

The reinforced concrete deck slab is designed per IRC~112:2011 (flexure, shear, crack width) and IRC~22:2014 (composite construction). Wheel loads are distributed using Pigeaud's method. The deck is checked for flexure in the transverse and longitudinal directions, punching shear, one-way (beam) shear, crack width, and reinforcement detailing.

\vspace{1em}
\begin{longtable}{|L{5.31cm}|p{9.66cm}|}%
\caption{\textbf{Deck Slab --- Loading and Geometry}}\\%
\hline%
\endfirsthead%
\caption[]{\textbf{Deck Slab --- Loading and Geometry} (continued)}\\%
\hline%
\endhead%
\multicolumn{2}{r}{{\footnotesize\itshape Continued on next page}} \\%
\endfoot%
\hline%
\endlastfoot%
\textnormal{Effective Span of Deck Slab, $l_{eff}$} & 2100 mm (girder spacing, c/c) \\*[6pt]
\hline
\textnormal{Deck Thickness, $t_s$} & 250.0 mm \\*[6pt]
\hline
\textnormal{Clear Cover (IRC 112 Cl. 15.2)} & Top 50 / Bottom 50 mm \\*[6pt]
\hline
\textnormal{Concrete Grade (IRC 112 Cl. 6.4)} & M40 ($f_{ck}$ = 40.0 MPa, $f_{ctm}$ = 3.0 MPa) \\*[6pt]
\hline
\textnormal{Reinforcement Grade (IRC 112 Cl. 6.2)} & Fe 500 ($f_y$ = 500 MPa) \\*[6pt]
\hline
\textnormal{Dead Load per Unit Area, $w_{DL}$} & 6.13 kN/m² (slab self-weight) \\*[6pt]
\hline
\textnormal{IRC 6 Wheel Load (Class A / 70R)} & 83.4 kN \\*[6pt]
\hline
\textnormal{Tyre Contact Width (IRC 6 Annex~A)} & 300 mm (transverse) \\*[6pt]
\hline
\textnormal{Impact Factor of the Governing Vehicle (IRC 6 Cl. 208)} & 1.100 \\*[6pt]
\hline
\textnormal{Governing Live Load Case} & Class70R(W) \\*[6pt]%
\end{longtable}

\vspace{1em}
\begin{longtable}{|C{2.64cm}|C{3.08cm}|C{2.64cm}|>{\centering\arraybackslash}p{3.70cm}|C{1.58cm}|}%
\caption{\textbf{Deck Slab --- Flexure Check: Interior Panel (Pigeaud's Method)}}\\%
\hline%
\rowcolor{tableHeaderBg}\hspace{0pt}\textbf{Location} & \hspace{0pt}\textbf{Parameter} & \hspace{0pt}\textbf{Formula / Reference} & \hspace{0pt}\textbf{Value} & \hspace{0pt}\textbf{Status} \\%
\hline%
\endfirsthead%
\caption[]{\textbf{Deck Slab --- Flexure Check: Interior Panel (Pigeaud's Method)} (continued)}\\%
\hline%
\rowcolor{tableHeaderBg}\hspace{0pt}\textbf{Location} & \hspace{0pt}\textbf{Parameter} & \hspace{0pt}\textbf{Formula / Reference} & \hspace{0pt}\textbf{Value} & \hspace{0pt}\textbf{Status} \\%
\hline%
\endhead%
\multicolumn{5}{r}{{\footnotesize\itshape Continued on next page}} \\%
\endfoot%
\hline%
\endlastfoot%
\multirow{5}{*}{\makecell{At Midspan\\(Sagging)}} & Transverse BM (DL), $M_{T,DL}$ & $w_{DL}\,l_{eff}^2/10$ & 2.70 kN-m/m & --- \\*[6pt]
\cline{2-5}\noalign{\nobreak}
 & Transverse BM (LL), $M_{T,LL}$ & Effective width (IRC 112 B3.1) & 25.63 kN-m/m & --- \\*[6pt]
\cline{2-5}\noalign{\nobreak}
 & Total Design BM, $M_{u,sag}$ & 1.35 DL + 1.50 LL & 45.94 kN-m/m & --- \\*[6pt]
\cline{2-5}\noalign{\nobreak}
 & Effective depth, $d$ & $t_s - c_{nom} - \phi/2$ & 194.0 mm & --- \\*[6pt]
\cline{2-5}\noalign{\nobreak}
 & Moment Capacity, $M_{Rd}$ & IRC 112 Cl. 12.2 & 52.23 kN-m/m & PASS \\*[6pt]
\hline
\multirow{3}{*}{\makecell{At Support\\(Hogging)}} & Total Design BM, $M_{u,hog}$ & 1.35 DL + 1.50 LL (at support) & 34.46 kN-m/m & --- \\*[6pt]
\cline{2-5}\noalign{\nobreak}
 & Required Top Steel, $A_{st,top}$ & $M_u / (0.87\,f_y\,d)$ & 488 mm²/m & --- \\*[6pt]
\cline{2-5}\noalign{\nobreak}
 & Moment Capacity, $M_{Rd}$ & IRC 112 Cl. 12.2 & 40.16 kN-m/m & PASS \\*[6pt]%
\end{longtable}
\noindent\textit{Note: IRC 112 Cl. 12.2. Distribution (longitudinal) reinforcement designed for 20\% of main steel moment (IRC 21 Cl. 305.18).}

\vspace{1em}
\begin{longtable}{|L{5.00cm}|C{3.18cm}|>{\centering\arraybackslash}p{4.09cm}|C{1.81cm}|}%
\caption{\textbf{Deck Slab --- Cantilever Overhang Flexure Check}}\\%
\hline%
\rowcolor{tableHeaderBg}\hspace{0pt}\textbf{Parameter} & \hspace{0pt}\textbf{Formula} & \hspace{0pt}\textbf{Value} & \hspace{0pt}\textbf{Status} \\%
\hline%
\endfirsthead%
\caption[]{\textbf{Deck Slab --- Cantilever Overhang Flexure Check} (continued)}\\%
\hline%
\rowcolor{tableHeaderBg}\hspace{0pt}\textbf{Parameter} & \hspace{0pt}\textbf{Formula} & \hspace{0pt}\textbf{Value} & \hspace{0pt}\textbf{Status} \\%
\hline%
\endhead%
\multicolumn{4}{r}{{\footnotesize\itshape Continued on next page}} \\%
\endfoot%
\hline%
\endlastfoot%
Overhang Length, $l_{oh}$ & --- & 1.05 m & --- \\*[6pt]
\hline
Crash Barrier Load Moment & IRC 6 Cl. 206.4 & 7.50 kN-m/m & --- \\*[6pt]
\hline
Dead Load Moment & $w_{DL}\,l_{oh}^2/2$ + railing & 4.92 kN-m/m & --- \\*[6pt]
\hline
Live Load Moment (eccentric wheel) & Wheel load $\times$ arm & 50.71 kN-m/m & --- \\*[6pt]
\hline
Total Hogging Moment, $M_{u,oh}$ & 1.35 DL + 1.50 (LL + CB) & 101.57 kN-m/m & --- \\*[6pt]
\hline
Moment Capacity (top steel), $M_{Rd,oh}$ & IRC 112 Cl. 12.2 & 108.27 kN-m/m & PASS \\*[6pt]%
\end{longtable}
\noindent\textit{Note: IRC 6 Cl. 206.4 crash barrier loads applied at kerb face; IRC 112 Cl. 12.2 flexure.}

\vspace{1em}
\begin{longtable}{|L{5.00cm}|C{3.18cm}|>{\centering\arraybackslash}p{4.09cm}|C{1.81cm}|}%
\caption{\textbf{Deck Slab --- Punching Shear Check (IRC~112 Cl.~10.4.6)}}\\%
\hline%
\rowcolor{tableHeaderBg}\hspace{0pt}\textbf{Parameter} & \hspace{0pt}\textbf{Formula / Reference} & \hspace{0pt}\textbf{Value} & \hspace{0pt}\textbf{Status} \\%
\hline%
\endfirsthead%
\caption[]{\textbf{Deck Slab --- Punching Shear Check (IRC~112 Cl.~10.4.6)} (continued)}\\%
\hline%
\rowcolor{tableHeaderBg}\hspace{0pt}\textbf{Parameter} & \hspace{0pt}\textbf{Formula / Reference} & \hspace{0pt}\textbf{Value} & \hspace{0pt}\textbf{Status} \\%
\hline%
\endhead%
\multicolumn{4}{r}{{\footnotesize\itshape Continued on next page}} \\%
\endfoot%
\hline%
\endlastfoot%
Design Wheel Load (ULS), $V_{Ed}$ & $\gamma_Q\,(1+IF)\,P_w$ & 137.6 kN & --- \\*[6pt]
\hline
Tyre Contact Area & $a \times b$ (IRC 6 Annex~A) & 300 $\times$ 150 mm & --- \\*[6pt]
\hline
Loaded Area at mid-depth, $b_0$ & $c_1 \times c_2$ (incl.\ WC dispersion) & 400 $\times$ 250 mm & --- \\*[6pt]
\hline
Control Perimeter, $u_1$ & $2(c_1+c_2) + 4\pi d$ & 3738 mm & --- \\*[6pt]
\hline
Punching Shear Stress, $v_{Ed}$ & $V_{Ed} / (u_1\,d)$ & 0.190 MPa & --- \\*[6pt]
\hline
Punching Resistance, $v_{Rd,c}$ & IRC 112 Eq.\ 10.1 & 0.555 MPa & --- \\*[6pt]
\hline
Punching Shear Check & $v_{Ed} \leq v_{Rd,c}$ & 0.34 & PASS \\*[6pt]%
\end{longtable}
\noindent\textit{Note: Punching shear reinforcement not typically required for deck slabs with $d \geq 200$ mm and adequate longitudinal reinforcement.}

\vspace{1em}
\clearpage
\begin{longtable}{|L{6.76cm}|>{\arraybackslash}p{8.21cm}|}%
\caption{\textbf{Crack Width Check (Deck Slab)}}\\%
\hline%
\rowcolor{tableHeaderBg}\hspace{0pt}\textbf{Parameter} & \hspace{0pt}\textbf{Value / Reference} \\%
\hline%
\endfirsthead%
\caption[]{\textbf{Crack Width Check (Deck Slab)} (continued)}\\%
\hline%
\rowcolor{tableHeaderBg}\hspace{0pt}\textbf{Parameter} & \hspace{0pt}\textbf{Value / Reference} \\%
\hline%
\endhead%
\multicolumn{2}{r}{{\footnotesize\itshape Continued on next page}} \\%
\endfoot%
\hline%
\endlastfoot%
\textnormal{Min. Reinforcement for Crack Control, $A_{s,min}$} & 303 mm²/m [IRC 112 Cl. 16.5.1] \\*[6pt]
\hline
\textnormal{Provided Reinforcement (bottom)} & $\phi$12 @ 175 mm c/c (646 mm²/m) \\*[6pt]
\hline
\textnormal{Max. Permissible Crack Width} & 0.30 mm \\*[6pt]
\hline
\textnormal{Calculated Crack Width, $w_k$ (governing)} & 0.2125 mm \\*[6pt]
\hline
\textnormal{Crack Width Check} & PASS \\*[6pt]%
\end{longtable}

\vspace{1em}
\begin{longtable}{|L{5.00cm}|C{3.18cm}|>{\centering\arraybackslash}p{4.09cm}|C{1.81cm}|}%
\caption{\textbf{One-Way (Beam) Shear Check (Deck Slab)}}\\%
\hline%
\rowcolor{tableHeaderBg}\hspace{0pt}\textbf{Parameter} & \hspace{0pt}\textbf{Formula / Reference} & \hspace{0pt}\textbf{Value} & \hspace{0pt}\textbf{Status} \\%
\hline%
\endfirsthead%
\caption[]{\textbf{One-Way (Beam) Shear Check (Deck Slab)} (continued)}\\%
\hline%
\rowcolor{tableHeaderBg}\hspace{0pt}\textbf{Parameter} & \hspace{0pt}\textbf{Formula / Reference} & \hspace{0pt}\textbf{Value} & \hspace{0pt}\textbf{Status} \\%
\hline%
\endhead%
\multicolumn{4}{r}{{\footnotesize\itshape Continued on next page}} \\%
\endfoot%
\hline%
\endlastfoot%
Design Shear per unit width, $V_{Ed}$ & $\gamma_{DL} V_{DL} + \gamma_{LL}(1{+}IF)V_{LL}$ & 89.24 kN/m & --- \\*[6pt]
\hline
Effective depth, $d$ & $t_s - c_{nom} - \phi/2$ & 194.0 mm & --- \\*[6pt]
\hline
Size factor, $k$ & $1 + \sqrt{200/d} \leq 2.0$ & 2.000 & --- \\*[6pt]
\hline
Long.\ reinforcement ratio, $\rho_l$ & $A_{sl}/(b_w\,d) \leq 0.02$ & 0.0033 & --- \\*[6pt]
\hline
Shear resistance (no stirrups), $V_{Rd,c}$ & $v_{Rd,c}\,b_w\,d$ (Cl.\ 10.3.2) & 107.58 kN/m & --- \\*[6pt]
\hline
One-Way Shear Check & $V_{Ed} \leq V_{Rd,c}$ & 0.83 & PASS \\*[6pt]%
\end{longtable}
\noindent\textit{Note: IRC 112 Cl. 10.3.2. Shear reinforcement not provided in deck slabs; capacity relies on concrete and main reinforcement.}

\vspace{1em}
\begin{longtable}{|L{5.00cm}|>{\centering\arraybackslash}p{3.72cm}|>{\centering\arraybackslash}p{3.72cm}|C{1.63cm}|}%
\caption{\textbf{Reinforcement Detailing Summary (Deck Slab)}}\\%
\hline%
\rowcolor{tableHeaderBg}\hspace{0pt}\textbf{Parameter} & \hspace{0pt}\textbf{Required / Limit} & \hspace{0pt}\textbf{Provided} & \hspace{0pt}\textbf{Status} \\%
\hline%
\endfirsthead%
\caption[]{\textbf{Reinforcement Detailing Summary (Deck Slab)} (continued)}\\%
\hline%
\rowcolor{tableHeaderBg}\hspace{0pt}\textbf{Parameter} & \hspace{0pt}\textbf{Required / Limit} & \hspace{0pt}\textbf{Provided} & \hspace{0pt}\textbf{Status} \\%
\hline%
\endhead%
\multicolumn{4}{r}{{\footnotesize\itshape Continued on next page}} \\%
\endfoot%
\hline%
\endlastfoot%
\multicolumn{4}{|l|}{\textbf{Main Reinforcement --- Bottom (Transverse)}} \\*[6pt]
\hline
Required Area, $A_{st,req}$ (mm²/m) & 596 mm²/m & 646 mm²/m & PASS \\*[6pt]
\hline
Bar Diameter $\times$ Spacing & $\phi \geq 10$ mm (IRC 112) & $\phi$12 @ 175 mm c/c & --- \\*[6pt]
\hline
Min.\ Reinforcement $A_{s,min}$ (IRC 112 Cl. 16.3.1) & 303 mm²/m & 646 mm²/m & PASS \\*[6pt]
\hline
Max.\ Bar Spacing (IRC 112 Cl. 16.3.2) & 300 mm & 175 mm & PASS \\*[6pt]
\hline
\multicolumn{4}{|l|}{\textbf{Distribution Reinforcement --- Longitudinal}} \\*[6pt]
\hline
Required Area, $A_{st,dist}$ (mm²/m) & $\geq 20\%$ of main steel & 377 mm²/m & PASS \\*[6pt]
\hline
\multicolumn{4}{|l|}{\textbf{Top Reinforcement (Support / Cantilever Overhang)}} \\*[6pt]
\hline
Required Area, $A_{st,top}$ (mm²/m) & 488 mm²/m & 492 mm²/m & PASS \\*[6pt]
\hline
\multicolumn{4}{|l|}{\textbf{Cover and Detailing}} \\*[6pt]
\hline
Clear Cover (IRC 112 Cl. 15.2) & $\geq$ 50 mm (Table 14.2) & Top 50 / Bottom 50 mm & PASS \\*[6pt]%
\end{longtable}
\noindent\textit{Note: IRC 112 Cl. 16.3, IS 456 Cl. 26.5. All reinforcement provisions satisfy strength and detailing requirements.}

% ===========================
\section{Cross Bracing Design}
\label{sec:cross-bracing}
% ===========================

Cross bracing between adjacent plate girders provides lateral stability during construction, resists transverse loads (wind, seismic, braking) in service, and prevents lateral torsional buckling of the girders. Members are designed per IS~800:2007 Cl.~7 (compression) and Cl.~6 (tension). Forces are derived from the grillage model under the governing load combination  (DL + LL + WL).

\vspace{1em}

\vspace{0.4em}
\noindent

\begin{longtable}{|C{2.0cm}|L{2.0cm}|L{2.2cm}|C{2.5cm}|C{2.0cm}|C{2.0cm}|}%
\caption{\textbf{Cross Bracing --- Connection and Section Properties}}\\%
\hline%
\rowcolor{tableHeaderBg}\hspace{0pt}\textbf{Panel} & \hspace{0pt}\textbf{Member} & \hspace{0pt}\textbf{Connection} & \hspace{0pt}\textbf{Section} & \hspace{0pt}\textbf{$A_g$ (mm²)} & \hspace{0pt}\textbf{$r_{min}$ (mm)} \\%
\hline%
\endfirsthead%
\caption[]{\textbf{Cross Bracing --- Connection and Section Properties} (continued)}\\%
\hline%
\rowcolor{tableHeaderBg}\hspace{0pt}\textbf{Panel} & \hspace{0pt}\textbf{Member} & \hspace{0pt}\textbf{Connection} & \hspace{0pt}\textbf{Section} & \hspace{0pt}\textbf{$A_g$ (mm²)} & \hspace{0pt}\textbf{$r_{min}$ (mm)} \\%
\hline%
\endhead%
\multicolumn{6}{r}{{\footnotesize\itshape Continued on next page}} \\%
\endfoot%
\endlastfoot%
\multirow{2}{*}{\makecell{G1-G2}} & Diagonal & Bolted & 40 x 40 x 3 & 237.0 & 8.0 \\*[6pt]\cline{2-6}\noalign{\nobreak}
\multirow{2}{*}{\makecell{G1-G2}} & Top / Bottom chord & Bolted & 40 x 40 x 3 & 237.0 & 8.0 \\*[6pt]\cline{2-6}
\hline\noalign{\nobreak}
\multirow{2}{*}{\makecell{G2-G3}} & Diagonal & Bolted & 40 x 40 x 3 & 237.0 & 8.0 \\*[6pt]\cline{2-6}\noalign{\nobreak}
\multirow{2}{*}{\makecell{G2-G3}} & Top / Bottom chord & Bolted & 40 x 40 x 3 & 237.0 & 8.0 \\*[6pt]\cline{2-6}
\hline\noalign{\nobreak}
\multirow{2}{*}{\makecell{G3-G4}} & Diagonal & Bolted & 40 x 40 x 3 & 237.0 & 8.0 \\*[6pt]\cline{2-6}\noalign{\nobreak}
\multirow{2}{*}{\makecell{G3-G4}} & Top / Bottom chord & Bolted & 40 x 40 x 3 & 237.0 & 8.0 \\*[6pt]\cline{2-6}
\hline\noalign{\nobreak}%
\end{longtable}
\noindent\textit{Note: $A_g$ = gross cross-sectional area; $r_{min}$ = minimum radius of gyration.}

\vspace{1em}
\begin{longtable}{|C{1.87cm}|L{1.87cm}|L{2.13cm}|C{2.13cm}|C{2.13cm}|>{\centering\arraybackslash}p{3.06cm}|}%
\caption{\textbf{Cross Bracing --- Slenderness Ratio Check (IS~800 Cl.~3.8 )}}\\%
\hline%
\rowcolor{tableHeaderBg}\hspace{0pt}\textbf{Panel} & \hspace{0pt}\textbf{Member} & \hspace{0pt}\textbf{Nature} & \hspace{0pt}\textbf{Eff.\ Length $KL$ (mm)} & \hspace{0pt}\textbf{$KL/r$} & \hspace{0pt}\textbf{Limit / Status} \\%
\hline%
\endfirsthead%
\caption[]{\textbf{Cross Bracing --- Slenderness Ratio Check (IS~800 Cl.~3.8 )} (continued)}\\%
\hline%
\rowcolor{tableHeaderBg}\hspace{0pt}\textbf{Panel} & \hspace{0pt}\textbf{Member} & \hspace{0pt}\textbf{Nature} & \hspace{0pt}\textbf{Eff.\ Length $KL$ (mm)} & \hspace{0pt}\textbf{$KL/r$} & \hspace{0pt}\textbf{Limit / Status} \\%
\hline%
\endhead%
\multicolumn{6}{r}{{\footnotesize\itshape Continued on next page}} \\%
\endfoot%
\hline%
\endlastfoot%
\multirow{3}{*}{\makecell{G1-G2}} & Diagonal & C & 2535 & 134.0 & 250 --- \textcolor{black}{PASS} \\*[6pt]
\cline{2-6}\noalign{\nobreak}
 & Top chord & C & 2100 & 111.0 & 250 --- \textcolor{black}{PASS} \\*[6pt]
\cline{2-6}\noalign{\nobreak}
 & Bottom chord & T & 2100 & 111.0 & 400 --- \textcolor{black}{PASS} \\*[6pt]
\hline
\multirow{3}{*}{\makecell{G2-G3}} & Diagonal & C & 2535 & 134.0 & 250 --- \textcolor{black}{PASS} \\*[6pt]
\cline{2-6}\noalign{\nobreak}
 & Top chord & C & 2100 & 111.0 & 250 --- \textcolor{black}{PASS} \\*[6pt]
\cline{2-6}\noalign{\nobreak}
 & Bottom chord & T & 2100 & 111.0 & 400 --- \textcolor{black}{PASS} \\*[6pt]
\hline
\multirow{3}{*}{\makecell{G3-G4}} & Diagonal & C & 2535 & 134.0 & 250 --- \textcolor{black}{PASS} \\*[6pt]
\cline{2-6}\noalign{\nobreak}
 & Top chord & C & 2100 & 111.0 & 250 --- \textcolor{black}{PASS} \\*[6pt]
\cline{2-6}\noalign{\nobreak}
 & Bottom chord & T & 2100 & 111.0 & 400 --- \textcolor{black}{PASS} \\*[6pt]%
\end{longtable}
\noindent\textit{Note:  3. Limit = 250 for compression members, 400 for tension members. $K = 1.0$ for members with both ends pinned.}

\vspace{1em}
\begin{longtable}{|C{1.58cm}|L{1.42cm}|C{1.73cm}|C{2.37cm}|C{1.42cm}|C{1.42cm}|C{0.94cm}|C{1.42cm}|}%
\caption{\textbf{Cross Bracing Design --- Capacity Summary}}\\%
\hline%
\rowcolor{tableHeaderBg}\hspace{0pt}\textbf{Panel} & \hspace{0pt}\textbf{Member} & \hspace{0pt}\textbf{Section} & \hspace{0pt}\textbf{Governing LC} & \hspace{0pt}\textbf{Demand (kN)} & \hspace{0pt}\textbf{Capacity (kN)} & \hspace{0pt}\textbf{UR} & \hspace{0pt}\textbf{Status} \\%
\hline%
\endfirsthead%
\caption[]{\textbf{Cross Bracing Design --- Capacity Summary} (continued)}\\%
\hline%
\rowcolor{tableHeaderBg}\hspace{0pt}\textbf{Panel} & \hspace{0pt}\textbf{Member} & \hspace{0pt}\textbf{Section} & \hspace{0pt}\textbf{Governing LC} & \hspace{0pt}\textbf{Demand (kN)} & \hspace{0pt}\textbf{Capacity (kN)} & \hspace{0pt}\textbf{UR} & \hspace{0pt}\textbf{Status} \\%
\hline%
\endhead%
\multicolumn{8}{r}{{\footnotesize\itshape Continued on next page}} \\%
\endfoot%
\hline%
\endlastfoot%
\multirow{2}{*}{\makecell{G1-G2}} & Diagonal & 40 x 40 x 3 & 1.5 EQ (b) & 8.049 & 16.810 & 0.740 & \textcolor{black}{PASS} \\*[6pt]\cline{2-8}\noalign{\nobreak}
 & Chord & 40 x 40 x 3 & 1.5 EQ (b) & 6.669 & 22.140 & 0.610 & \textcolor{black}{PASS} \\[6pt]\cline{2-8}
\hline
\multirow{2}{*}{\makecell{G2-G3}} & Diagonal & 40 x 40 x 3 & 1.5 EQ (a) & 7.159 & 45.000 & 0.160 & \textcolor{black}{PASS} \\*[6pt]\cline{2-8}\noalign{\nobreak}
 & Chord & 40 x 40 x 3 & 1.5 EQ (a) & 5.931 & 45.000 & 0.130 & \textcolor{black}{PASS} \\[6pt]\cline{2-8}
\hline
\multirow{2}{*}{\makecell{G3-G4}} & Diagonal & 40 x 40 x 3 & 1.5 EQ (b) & 9.644 & 45.000 & 0.210 & \textcolor{black}{PASS} \\*[6pt]\cline{2-8}\noalign{\nobreak}
 & Chord & 40 x 40 x 3 & 1.5 EQ (b) & 7.990 & 45.000 & 0.180 & \textcolor{black}{PASS} \\[6pt]\cline{2-8}%
\end{longtable}
\noindent\textit{Note: Designed per IS 800 Cl. 7 (compression) and Cl. 6 (tension). OsdagBridge cross-bracing module used.}

% ===========================
\section{End Diaphragm Design}
\label{sec:end-diaphragm}
% ===========================

End diaphragms at the supports transfer transverse loads to the bearings, restrain the bottom flanges against lateral displacement, and maintain the girder cross-section geometry during construction and in service. They are designed per IS~800:2007 and IRC~24:2010 Cl.~507.

\vspace{1em}
\begin{longtable}{|C{1.20cm}|C{1.12cm}|L{2.58cm}|L{3.97cm}|C{1.29cm}|C{1.29cm}|C{1.29cm}|}%
\caption{\textbf{End Diaphragm --- Members and Section Properties (input)}}\\%
\hline%
\rowcolor{tableHeaderBg}\hspace{0pt}\textbf{Pair} & \hspace{0pt}\textbf{Member} & \hspace{0pt}\textbf{Type} & \hspace{0pt}\textbf{Section(s)} & \hspace{0pt}\textbf{$A$ (cm\textsuperscript{2})} & \hspace{0pt}\textbf{$r_y$ (cm)} & \hspace{0pt}\textbf{$r_z$ (cm)} \\%
\hline%
\endfirsthead%
\caption[]{\textbf{End Diaphragm --- Members and Section Properties (input)} (continued)}\\%
\hline%
\rowcolor{tableHeaderBg}\hspace{0pt}\textbf{Pair} & \hspace{0pt}\textbf{Member} & \hspace{0pt}\textbf{Type} & \hspace{0pt}\textbf{Section(s)} & \hspace{0pt}\textbf{$A$ (cm\textsuperscript{2})} & \hspace{0pt}\textbf{$r_y$ (cm)} & \hspace{0pt}\textbf{$r_z$ (cm)} \\%
\hline%
\endhead%
\multicolumn{7}{r}{{\footnotesize\itshape Continued on next page}} \\%
\endfoot%
\hline%
\endlastfoot%
G1G2 & E1M1 & Cross Bracing (X-Bracing) & bracing:\allowbreak{} IS 100 x 100 x 10; top chord:\allowbreak{} IS 100 x 100 x 10; bottom chord:\allowbreak{} IS 100 x 100 x 10 & 19.10 & 3.07 & 3.07 \\*[6pt]
\hline
G1G2 & E1M2 & Cross Bracing (X-Bracing) & bracing:\allowbreak{} IS 100 x 100 x 10; top chord:\allowbreak{} IS 100 x 100 x 10; bottom chord:\allowbreak{} IS 100 x 100 x 10 & 19.10 & 3.07 & 3.07 \\*[6pt]
\hline
G2G3 & E2M1 & Cross Bracing (X-Bracing) & bracing:\allowbreak{} IS 100 x 100 x 10; top chord:\allowbreak{} IS 100 x 100 x 10; bottom chord:\allowbreak{} IS 100 x 100 x 10 & 19.10 & 3.07 & 3.07 \\*[6pt]
\hline
G2G3 & E2M2 & Cross Bracing (X-Bracing) & bracing:\allowbreak{} IS 100 x 100 x 10; top chord:\allowbreak{} IS 100 x 100 x 10; bottom chord:\allowbreak{} IS 100 x 100 x 10 & 19.10 & 3.07 & 3.07 \\*[6pt]
\hline
G3G4 & E3M1 & Cross Bracing (X-Bracing) & bracing:\allowbreak{} IS 100 x 100 x 10; top chord:\allowbreak{} IS 100 x 100 x 10; bottom chord:\allowbreak{} IS 100 x 100 x 10 & 19.10 & 3.07 & 3.07 \\*[6pt]
\hline
G3G4 & E3M2 & Cross Bracing (X-Bracing) & bracing:\allowbreak{} IS 100 x 100 x 10; top chord:\allowbreak{} IS 100 x 100 x 10; bottom chord:\allowbreak{} IS 100 x 100 x 10 & 19.10 & 3.07 & 3.07 \\*[6pt]%
\end{longtable}
\noindent\textit{Note: section properties as stored with the input (section database values).}

\vspace{1em}
\begin{longtable}{|C{2.04cm}|C{1.44cm}|C{1.70cm}|C{2.21cm}|C{1.36cm}|C{1.44cm}|C{0.93cm}|C{1.19cm}|}%
\caption{\textbf{End Diaphragm Design --- Capacity Summary}}\\%
\hline%
\rowcolor{tableHeaderBg}\hspace{0pt}\textbf{Pair / Member} & \hspace{0pt}\textbf{Force} & \hspace{0pt}\textbf{Section} & \hspace{0pt}\textbf{Governing LC} & \hspace{0pt}\textbf{Demand (kN)} & \hspace{0pt}\textbf{Capacity (kN)} & \hspace{0pt}\textbf{UR} & \hspace{0pt}\textbf{Status} \\%
\hline%
\endfirsthead%
\caption[]{\textbf{End Diaphragm Design --- Capacity Summary} (continued)}\\%
\hline%
\rowcolor{tableHeaderBg}\hspace{0pt}\textbf{Pair / Member} & \hspace{0pt}\textbf{Force} & \hspace{0pt}\textbf{Section} & \hspace{0pt}\textbf{Governing LC} & \hspace{0pt}\textbf{Demand (kN)} & \hspace{0pt}\textbf{Capacity (kN)} & \hspace{0pt}\textbf{UR} & \hspace{0pt}\textbf{Status} \\%
\hline%
\endhead%
\multicolumn{8}{r}{{\footnotesize\itshape Continued on next page}} \\%
\endfoot%
\hline%
\endlastfoot%
\multicolumn{8}{|p{15.2cm}|}{\textit{No end-diaphragm design results: The analysis returned no end-diaphragm member forces, so no end-diaphragm Demand/Capacity checks were performed.}} \\*[6pt]%
\end{longtable}
\noindent\textit{Note: end-diaphragm results come from \texttt{end\_diaphragm\_design\_results} only; cross-bracing results are not reused here (see Section 5.3).}


% The overall summary starts on a new page (its heading was left alone at a
% page bottom; see the note above Section 5.1 on why \clearpage is used).
\clearpage
% ===========================
\section{Overall Design Check Summary}
\label{sec:overall-summary}
% ===========================

\vspace{1em}
\begin{longtable}{|C{3.29cm}|L{4.35cm}|C{2.22cm}|C{2.22cm}|>{\centering\arraybackslash}p{1.54cm}|}%
\caption{\textbf{Overall Design Check Summary --- All Members}}\label{tab:design-check-summary}\\%
\hline%
\rowcolor{tableHeaderBg}\hspace{0pt}\textbf{Member / Check} & \hspace{0pt}\textbf{Governing Load Combo} & \hspace{0pt}\textbf{Demand} & \hspace{0pt}\textbf{Capacity} & \hspace{0pt}\textbf{UR} \\%
\hline%
\endfirsthead%
\caption[]{\textbf{Overall Design Check Summary --- All Members} (continued)}\\%
\hline%
\rowcolor{tableHeaderBg}\hspace{0pt}\textbf{Member / Check} & \hspace{0pt}\textbf{Governing Load Combo} & \hspace{0pt}\textbf{Demand} & \hspace{0pt}\textbf{Capacity} & \hspace{0pt}\textbf{UR} \\%
\hline%
\endhead%
\multicolumn{5}{r}{{\footnotesize\itshape Continued on next page}} \\%
\endfoot%
\hline%
\endlastfoot%
Girder --- Moment & ACCIDENTAL\_\allowbreak{}1:\allowbreak{} 1.0DL + 1.0DW + 0.75LL & 2224.95 kNm & 11371.85 kNm & 0.20 \\*[6pt]
\hline
Girder --- Shear & ACCIDENTAL\_\allowbreak{}1:\allowbreak{} 1.0DL + 1.0DW + 0.75LL & 379.98 kN & 720.21 kN & 0.53 \\*[6pt]
\hline
Girder --- LTB (constr.) & SW & 206.16 kNm & 900.22 kNm & 0.23 \\*[6pt]
\hline
Girder --- Deflection & 1.0 DL + 1.0 LL & 22.87 mm & 50.00 mm & 0.46 \\*[6pt]
\hline
Girder --- Stress & SLS\_\allowbreak{}RARE\_\allowbreak{}1:\allowbreak{} 1.0DL + 1.2DW + 1.0LL + 0.6WL & 81.15 MPa & 315.00 MPa & 0.26 \\*[6pt]
\hline
Girder --- Fatigue & SLS\_\allowbreak{}FREQUENT\_\allowbreak{}1:\allowbreak{} 1.0DL + 1.2DW + 0.75LL + 0.5WL & 145.86 MPa & 92.49 MPa & \textcolor{red}{1.58} \\*[6pt]
\hline
Transverse Shear (slab) & --- & --- & --- & --- \\*[6pt]
\hline
Crack Width (slab) & SLS\_\allowbreak{}FREQUENT\_\allowbreak{}1:\allowbreak{} 1.0DL + 1.2DW + 0.75LL + 0.5WL & 0.212 mm & 0.300 mm & 0.71 \\*[6pt]
\hline
Deck --- Flexure (sagging) & Basic ULS: 1.35DL + 1.5LL & 45.94 kN-m/m & 52.23 kN-m/m & 0.88 \\*[6pt]
\hline
Deck --- Flexure (hogging) & Basic ULS: 1.35DL + 1.5LL & 34.46 kN-m/m & 40.16 kN-m/m & 0.86 \\*[6pt]
\hline
Deck --- Cantilever Overhang & Basic ULS: 1.35DL + 1.5LL & 101.57 kN-m/m & 108.27 kN-m/m & 0.94 \\*[6pt]
\hline
Deck --- Punching Shear & Basic ULS: 1.35DL + 1.5LL & 0.19 MPa & 0.55 MPa & 0.34 \\*[6pt]
\hline
Deck --- One-Way Shear & Basic ULS: 1.35DL + 1.5LL & 89.24 kN/m & 107.58 kN/m & 0.83 \\*[6pt]
\hline
Cross Bracing --- Compression & 1.5 EQ (b) & 12.44 kN & 16.810 kN & 0.74 \\*[6pt]
\hline
Cross Bracing --- Tension & 1.5 EQ (b) & 9.45 kN & 45.000 kN & 0.21 \\*[6pt]
\hline
Cross Bracing --- Slenderness & --- & 134.0 & 250 & 0.54 \\*[6pt]
\hline
End Diaphragm & \multicolumn{4}{c|}{\parbox[t]{9.2cm}{\centering\small The analysis returned no end-diaphragm member forces, so no end-diaphragm Demand/Capacity checks were performed.\strut}} \\*[6pt]%
\end{longtable}
\noindent\textit{Note: UR = Demand / Capacity. All values $\leq 1.0$ indicate passing checks. The governing check for each component is highlighted in the individual design check sections above.}

\subsection*{Utilization Ratio Charts}
\addcontentsline{toc}{subsection}{Utilization Ratio Charts}
The charts below plot the Utilization Ratio (UR = Demand / Capacity) of every design check of the four primary structural elements. Green bars pass (UR $\leq$ 1.0), red hatched bars fail (UR $>$ 1.0) and the red dashed line marks the limit UR = 1.0. Girder bars show the governing girder of each check (in brackets). Checks without valid Demand / Capacity data are shown as \emph{no data} and are never counted as passing.

\begin{figure}[H]
\centering
\includegraphics[width=0.95\textwidth,keepaspectratio]{C:/Users/udayk/AppData/Local/Temp/tmpymqxq2pd/images/charts/ur_girders.png}
\caption{Utilization Ratio --- Steel Plate Girders}
\end{figure}

\begin{figure}[H]
\centering
\includegraphics[width=0.95\textwidth,keepaspectratio]{C:/Users/udayk/AppData/Local/Temp/tmpymqxq2pd/images/charts/ur_deck.png}
\caption{Utilization Ratio --- Concrete Deck Slab}
\end{figure}

\begin{figure}[H]
\centering
\includegraphics[width=0.95\textwidth,keepaspectratio]{C:/Users/udayk/AppData/Local/Temp/tmpymqxq2pd/images/charts/ur_cross_bracing.png}
\caption{Utilization Ratio --- Cross Bracing}
\end{figure}

\begin{figure}[H]
\centering
\includegraphics[width=0.95\textwidth,keepaspectratio]{C:/Users/udayk/AppData/Local/Temp/tmpymqxq2pd/images/charts/ur_end_diaphragms.png}
\caption{Utilization Ratio --- End Diaphragms}\par\vspace{2pt}{\footnotesize\itshape The analysis returned no end-diaphragm member forces, so no end-diaphragm Demand/Capacity checks were performed.}
\end{figure}


\chapter{Drawings and Visualizations}
\label{ch:drawings}

This section presents CAD-generated views of the designed bridge and its components. All views are generated automatically by OsdagBridge using pythonOCC.

\section{Bridge Configuration and Layout}
\label{sec:bridge-layout}

\begin{figure}[H]
\centering
\vspace{4pt}
\includegraphics[width=0.85\textwidth]{C:/Users/udayk/AppData/Local/Temp/tmpymqxq2pd/images/final_geometry.png}
\caption{Overall 3D Bridge Superstructure}
\label{fig:bridge-3d}
\end{figure}


\begin{figure}[H]
\centering
\vspace{4pt}
\includegraphics[width=0.85\textwidth]{C:/Users/udayk/AppData/Local/Temp/tmpymqxq2pd/images/cross_section.png}
\caption{Typical Cross Section}
\label{fig:cross-section}
\end{figure}


\begin{figure}[H]
\centering
\vspace{4pt}
\includegraphics[width=0.85\textwidth]{C:/Users/udayk/AppData/Local/Temp/tmpymqxq2pd/images/girder_top.png}
\caption{Top View}
\label{fig:top-view}
\end{figure}


\section{Plate Girder --- Detailed Views}
\label{sec:girder-views}

\begin{figure}[H]
\centering
\vspace{4pt}
\includegraphics[width=0.85\textwidth]{C:/Users/udayk/AppData/Local/Temp/tmpymqxq2pd/images/girder_3d.png}
\caption{3D View of Plate Girders}
\label{fig:girder-3d}
\end{figure}


\begin{figure}[H]
\centering
\vspace{4pt}
\includegraphics[width=0.85\textwidth]{C:/Users/udayk/AppData/Local/Temp/tmpymqxq2pd/images/section_preview.png}
\caption{Cross Section of Plate Girder}
\label{fig:girder-xsec}
\end{figure}


\begin{figure}[H]
\centering
\vspace{4pt}
\includegraphics[width=0.85\textwidth]{C:/Users/udayk/AppData/Local/Temp/tmpymqxq2pd/images/stiffener_preview.png}
\caption{Side View of Girder}
\label{fig:girder-side}
\end{figure}


\section{Cross Bracing Detail}
\label{sec:bracing-detail}

\begin{figure}[H]
\centering
\vspace{4pt}
\includegraphics[width=0.85\textwidth]{C:/Users/udayk/AppData/Local/Temp/tmpymqxq2pd/images/cb_diagram.png}
\caption{Cross Bracing Layout}
\label{fig:cb-layout}
\end{figure}


\noindent
\begin{minipage}[t]{0.31\textwidth}
\centering
\includegraphics[width=\linewidth]{C:/Users/udayk/AppData/Local/Temp/tmpymqxq2pd/images/cb_bracing.png}
\captionof{figure}{Cross Bracing -- Bracing Section}
\label{fig:cb-bracing-sec}
\end{minipage}
\hfill
\begin{minipage}[t]{0.31\textwidth}
\centering
\includegraphics[width=\linewidth]{C:/Users/udayk/AppData/Local/Temp/tmpymqxq2pd/images/cb_top_chord.png}
\captionof{figure}{Cross Bracing -- Top Chord Section}
\label{fig:cb-top-chord}
\end{minipage}
\hfill
\begin{minipage}[t]{0.31\textwidth}
\centering
\includegraphics[width=\linewidth]{C:/Users/udayk/AppData/Local/Temp/tmpymqxq2pd/images/cb_bottom_chord.png}
\captionof{figure}{Cross Bracing -- Bottom Chord Section}
\label{fig:cb-bottom-chord}
\end{minipage}
            

\section{End Diaphragm Detail}
\label{sec:diaphragm-detail}

\begin{figure}[H]
\centering
\vspace{4pt}
\includegraphics[width=0.85\textwidth]{C:/Users/udayk/AppData/Local/Temp/tmpymqxq2pd/images/ed_diagram.png}
\caption{End Diaphragm Layout}
\label{fig:ed-layout}
\end{figure}


\noindent
\begin{minipage}[t]{0.31\textwidth}
\centering
\includegraphics[width=\linewidth]{C:/Users/udayk/AppData/Local/Temp/tmpymqxq2pd/images/ed_bracing.png}
\captionof{figure}{End Diaphragm -- Bracing Section}
\label{fig:ed-bracing-sec}
\end{minipage}
\hfill
\begin{minipage}[t]{0.31\textwidth}
\centering
\includegraphics[width=\linewidth]{C:/Users/udayk/AppData/Local/Temp/tmpymqxq2pd/images/ed_top_chord.png}
\captionof{figure}{End Diaphragm -- Top Chord Section}
\label{fig:ed-top-chord}
\end{minipage}
\hfill
\begin{minipage}[t]{0.31\textwidth}
\centering
\includegraphics[width=\linewidth]{C:/Users/udayk/AppData/Local/Temp/tmpymqxq2pd/images/ed_bottom_chord.png}
\captionof{figure}{End Diaphragm -- Bottom Chord Section}
\label{fig:ed-bottom-chord}
\end{minipage}



\chapter{Material Take-off \& Quantity Summary}
\label{ch:material-takeoff}

\begingroup
\setlength{\tabcolsep}{3.5pt}
\begin{longtable}{|C{0.95cm}|L{3.64cm}|C{2.49cm}|C{1.72cm}|C{1.72cm}|C{1.72cm}|C{1.72cm}|}%
\caption{\textbf{Bill of Materials (Steel, Concrete, and Reinforcement Quantities)}}\label{tab:bill-of-materials}\\%
\hline%
\rowcolor{tableHeaderBg}\hspace{0pt}\textbf{S.N.} & \hspace{0pt}\textbf{Item Description} & \hspace{0pt}\textbf{Volume} & \hspace{0pt}\textbf{Quantity} & \hspace{0pt}\textbf{Total Volume} & \hspace{0pt}\textbf{Weight (MT)} & \hspace{0pt}\textbf{Total Weight (MT)} \\%
\hline%
\endfirsthead%
\caption[]{\textbf{Bill of Materials (Steel, Concrete, and Reinforcement Quantities)} (continued)}\\%
\hline%
\rowcolor{tableHeaderBg}\hspace{0pt}\textbf{S.N.} & \hspace{0pt}\textbf{Item Description} & \hspace{0pt}\textbf{Volume} & \hspace{0pt}\textbf{Quantity} & \hspace{0pt}\textbf{Total Volume} & \hspace{0pt}\textbf{Weight (MT)} & \hspace{0pt}\textbf{Total Weight (MT)} \\%
\hline%
\endhead%
\multicolumn{7}{r}{{\footnotesize\itshape Continued on next page}} \\%
\endfoot%
\hline%
\endlastfoot%
1 & Structural Steel (IS 2062) for Girders & $0.03870\text{ m}^2 \times 30.00\text{ m} = 1.16100\text{ m}^3$ & 4 & 4.64 & 9.11 & 36.46 \\*
\hline
2(a) & Cross Bracing - Top Chord & $0.00024\text{ m}^2 \times 2.10\text{ m} = 0.00050\text{ m}^3$ & 27 & 0.01 & 0.0039 & 0.11 \\*
\hline
2(b) & Cross Bracing - Bottom Chord & $0.00024\text{ m}^2 \times 2.10\text{ m} = 0.00050\text{ m}^3$ & 27 & 0.01 & 0.0039 & 0.11 \\*
\hline
2(c) & Cross Bracing - Diagonal Chord & $0.00024\text{ m}^2 \times 2.53\text{ m} = 0.00060\text{ m}^3$ & 54 & 0.03 & 0.0047 & 0.25 \\*
\hline
3 & End Diaphragms & N.A. & N.A. & N.A. & N.A. & N.A. \\*
\hline
4 & Concrete (M40) for Deck Slab & $8.40\text{ m} \times 0.25\text{ m} \times 30.00\text{ m} = 63.00\text{ m}^3$ & 1 & 63.00 & 157.50 & 157.50 \\*
\hline
5 & Reinforcement Steel (Fe 500) & $0.032102\text{ m}^2 \times 30.00\text{ m} = 0.96306\text{ m}^3$ & 1 & 0.96 & 7.56 & 7.56 \\*
\hline
6 & Shear Stud Connectors & $0.000314\text{ m}^2 \times 0.100\text{ m} = 0.000031\text{ m}^3$ & 640 & 0.02 & 0.000247 & 0.158 \\*
\hline
7 & Crash Barrier & $0.30230\text{ m}^2 \times 30.00\text{ m} = 9.06900\text{ m}^3$ & 2 & 18.14 & 22.67 & 45.35 \\*%
\end{longtable}
\endgroup
\section{Material Quantity Charts}
The charts below are generated from the quantities in the Bill of Materials above. Each material is shown in its own unit (structural and reinforcement steel in metric tonnes, concrete in m\textsuperscript{3}); quantities that could not be computed are marked \emph{not available}.

\begin{figure}[H]
\centering
\includegraphics[width=0.80\textwidth,keepaspectratio]{C:/Users/udayk/AppData/Local/Temp/tmpymqxq2pd/images/charts/qty_1.png}
\caption{Structural Steel Tonnage (MT)}\par\vspace{2pt}{\footnotesize\itshape End-diaphragm tonnage not available (N.A. in Table 7.1):\allowbreak{} the member lengths of the cross-bracing end-diaphragm frames are not in the design output.}
\end{figure}

\begin{figure}[H]
\centering
\includegraphics[width=0.80\textwidth,keepaspectratio]{C:/Users/udayk/AppData/Local/Temp/tmpymqxq2pd/images/charts/qty_2.png}
\caption{Concrete Volume (m³)}
\end{figure}

\begin{figure}[H]
\centering
\includegraphics[width=0.80\textwidth,keepaspectratio]{C:/Users/udayk/AppData/Local/Temp/tmpymqxq2pd/images/charts/qty_3.png}
\caption{Reinforcement Steel (MT)}\par\vspace{2pt}{\footnotesize\itshape Estimated by the BOQ module at 120 kg of reinforcement per m³ of deck concrete.}
\end{figure}


\chapter{Standards \& Assumptions}
\label{ch:Design Standards}

This section provides references to standards used to calibrate the OsdagBridge
design modules, and notes the limitations of the current software version.

\section{Design Standards}
\label{sec:design_standards}

The following Indian Road Congress (IRC) codes and Indian Standards (IS) 
form the basis of all design calculations in this software.

\vspace{0.5cm}

\begingroup
\setlength{\tabcolsep}{3.5pt}
\begin{table}[H]
\caption{\textbf{IRC Codes}}
\begin{tabular}{|L{2.9cm}|C{1.2cm}|L{10.9cm}|}
\hline
\textbf{Code} & \textbf{Year} & \textbf{Title / Scope} \\ 
\hline
IRC 5 & 2015 & General Features of Design - carriageway widths, kerb, footpath dimensions \\ 
\hline
IRC 6 & 2017 & Loads and Load Combinations - dead load, live load, impact, wind, temperature, etc. \\ 
\hline
IRC 22 & 2015 & Composite Construction (LS) - Composite section properties, ULS/SLS design, shear connectors \\ 
\hline
IRC 24 & 2010 & Steel Road Bridges (LS) - Stiffener design, skew angle limits, diaphragm requirements \\ 
\hline
IRC 112 & 2020 & Concrete Road Bridges - Deck slab flexure, shear, crack width, reinforcement \\ 
\hline
IRC SP 114 & 2018 & Seismic Design of Road Bridges \\ 
\hline
\end{tabular}
\end{table}

\begingroup
\setlength{\tabcolsep}{3.5pt}
\begin{table}[H]
\caption{\textbf{IS Codes}}
\begin{tabular}{|L{2.9cm}|C{1.2cm}|L{10.9cm}|}
\hline
\textbf{Code} & \textbf{Year} & \textbf{Scope} \\
\hline
IS 800 & 2007 & Steel construction - tension, compression, bending, shear, LTB, stiffeners, combined checks \\
\hline
IS 456 & 2000 & Concrete - simplified stress-block for deck moment capacity \\
\hline
IS 1786 & 2008 & Reinforcement steel properties \\
\hline
IS 1893 (Part 3) & 2014 & Earthquake resistant design \\
\hline
IS 2062 & 2011 & Structural steel - yield and ultimate strength by grade \\
\hline
\end{tabular}
\end{table}

\clearpage
\section{Analysis and Design Assumptions of This Version}
\label{sec:assumptions}

\textbf{Structural Analysis}

\begin{itemize}
    \item All girders are modelled as simply supported; continuous spans are not currently supported.
    \item A 3D grillage model (OSPGrillage) is used for load distribution. Grillage members carry composite section properties after the construction stage.
    \item The transverse member forces are computed using an approximate 2D frame analogy. For irregular or skewed geometries, a 3D FEM is recommended.
    \item Fixed bearing stiffness is modelled as $k = 1{,}000{,}000 \,\text{kN/m}$ (virtually rigid); free/expansion bearing as $k = 100 \, \text{kN/m}$.
    \item Construction stage sequence analysis is approximate. Detailed staged analysis should be performed for long-term deflection checks.
\end{itemize}

\textbf{Composite Action}

\begin{itemize}
    \item Full shear connection is assumed at ULS with headed stud connectors designed per IRC 22:2015 Cl.606.
    \item Short-term composite section properties (modular ratio $n = E_s/E_{cm}$) are used for ULS checks.
    \item Long-term composite section properties (with creep-adjusted modular ratio) are used for SLS deflection and crack-width checks.
    \item Pre-composite stage: the steel girder alone resists all construction loads prior to concrete gaining strength.
\end{itemize}


\textbf{Material Properties (IRC 22:2015 Annex III)}

\begin{itemize}
\item Steel: $E_s = 200{,}000 \, \text{MPa}, \ G_s = 80{,}000 \, \text{MPa}, \ \nu = 0.30, \ \alpha = 11.7 \times 10^{-6}/^\circ\mathrm{C}$ (grade-independent).
\item Minimum structural concrete grade: M25.
\item Default reinforcement grade: Fe500.
\end{itemize}

\textbf{Partial Safety Factors (IRC 22:2015 Cl.601.4)}

\begin{itemize}
\item $\gamma_{m0}$ (steel yield, ULS) = 1.10
\item $\gamma_{m1}$ (steel ultimate, ULS) = 1.25
\item Reinforcement (ULS) = 1.15
\item Welds -- shop: 1.25; field: 1.50
\item Fatigue ($\gamma_{mft}$) = 1.35
\end{itemize}

\textbf{Loading}

\begin{itemize}
\item Dead load densities per IRC 6:2017 Cl.203: structural steel $78.5 \ \text{kN/m}^3$, concrete $25.0 \ \text{kN/m}^3$, bituminous wearing course $24.0 \ \text{kN/m}^3$.
\item Multi-lane live load reduction factors per IRC 6:2017 Cl.204.4 Table 6A: 1st lane = 1.0, 2nd lane = 0.8, 3rd lane onwards = 0.4.
\item Impact factor (dynamic load allowance) computed from span per IRC 6:2017 Cl.208.2/208.3.
\item Wind load applied as transverse, longitudinal, and vertical components per IRC 6:2017 Cl.209.3.3--209.3.5.
\end{itemize}

\textbf{Serviceability Limits}

\begin{itemize}
\item Deflection limits (IRC 22:2015 Cl.604.3.2): live load + impact $\leq L/800$; total $\leq L/600$.
\item SLS stress limits: concrete $\sigma_c \leq 0.48 f_{ck}$; rebar $\sigma_s \leq 0.80 f_{yk}$; steel $f_e \leq 0.9 f_y$.
\item Permissible crack width: $w_k \leq 0.3 \ \text{mm}$ (bridge deck, exposure class XS2/XD2 per IRC 112:2020 Cl.12.3.2).
\end{itemize}

\textbf{Fatigue}

\begin{itemize}
\item Reference fatigue life: $N_{sc} = 2 \times 10^6$ cycles (IRC 22:2015 Cl.605).
\item Constant stress range is assumed. A thickness correction factor $\mu_r$ is applied for plate thickness $> 25 \ \text{mm}$.
\end{itemize}

\textbf{Stiffener Design}

\begin{itemize}
\item Intermediate transverse stiffener and bearing stiffener design follows IS 800:2007 Cl.8.7.2/8.7.3 and IRC 24:2010 Cl.509.7.2/509.7.3.
\end{itemize}


\section{Known Limitations of This Version}
\label{sec:limitations}

\begin{itemize}
\item Substructure (piers, pile caps, foundations) and bearing design are not included.
\item Splice connection design is not implemented.
\item Skew angle $>$ 15 degrees requires independent manual analysis (IRC 24 Cl. 504.8).
\item Construction stage sequence analysis is approximate; detailed staged analysis
  should be performed for long-term deflection checks.
\item The grillage analysis assumes simply supported boundary conditions;
  continuous spans are not currently supported.
\end{itemize}


\chapter*{References}
\addcontentsline{toc}{chapter}{References}


\begin{enumerate}

\item IRC 5 (2024) --- \textit{Standard Specifications and Code of Practice for Road Bridges, Section I: General Features of Design.}

\item IRC 6 (2017) --- \textit{Standard Specifications and Code of Practice for Road Bridges, Section II: Loads and Load Combinations.}

\item IRC 22 (2014) --- \textit{Standard Specifications and Code of Practice for Road Bridges, Section VI: Composite Construction (Limit State Design).}

\item IRC 24 (2010) --- \textit{Standard Specifications and Code of Practice for Road Bridges, Section V: Steel Road Bridges (Limit State Method).}

\item IRC 112 (2011) --- \textit{Code of Practice for Concrete Road Bridges.}

\item IRC SP 114 (2018) --- \textit{Guidelines for Seismic Design of Road Bridges.}

\item IS 800 (2007) --- \textit{Indian Standard: General Construction in Steel --- Code of Practice.}

\item IS 2062 (2011) --- \textit{Hot Rolled Medium and High Tensile Structural Steel --- Specification.}

\item Subramanian, N. (2008). \textit{Design of Steel Structures.} Oxford University Press.

\item Steel-INSDAG Teaching Resource Materials. \url{https://www.steel-insdag.org}

\end{enumerate}

\end{document}